\chapter{応用地質・地盤}
\section{岩石の成因を地盤の挙動へつなぐ}
応用地質では、材料の強さだけでなく割れ目、弱層、水、応力の向きなどが構造物や斜面の安定性にどう影響するかを調べる。硬い鉱物からできた岩石でも、多数の割れ目を持つ岩盤では弱面に沿う破壊が支配的になり得る。地質名をそのまま一つの強度値と同一視しない。

\section{特殊土の名称を母材と堆積過程に戻す}
\label{theory:engineering:soil}
まさ土は花崗岩類の風化で粒状化した土である\cite{mlit_masa}。しらすは主に未溶結の火砕流堆積物などの火山性物質であり、花崗岩風化とは違う。しらすの再移動・再堆積もある\cite{maff_shiras2017,gsj_shiras1957}。白い砂状だから同じ材料という判断を避ける。
黄土（レス）は風成の細粒堆積物である\cite{gsj_loess2002,loess2003}。火山灰由来の粘性土と同義ではない。泥炭（ピート）は植物遺体が十分に分解されず蓄積した有機質土であり、植物繊維を残すことがある\cite{mlit_peat}。分解程度の幅もある。残留土は生成場所に残り、崩積土は斜面を移動して集積した土であり、名称が示す母材と後の移動を別に整理する。

\section{風化・熱水変質と強度}
\label{theory:engineering:weathering}
風化は地表近くの水・酸素・温度変化などの環境に岩石が応答する過程であり、熱水変質は温かい流体との反応を中心とする。どちらも鉱物の変化や結合の変化を起こす。粘土化・割れ目の増加などで弱化しやすい一方、珪化などで硬化する場合もある\cite{jseg_engineering_faq}。
有色鉱物の多い岩石が風化に不安定という一般的比較は、鉱物の形成条件と地表環境の違いを反映する。実際の耐風化性には粒子の結合、割れ目、流体経路が関係するので色だけで寿命を決めない。

\section{スレーキング：空気と固体の応力を分ける}
\label{theory:engineering:slaking}
乾燥した岩塊を水に浸すと崩れるスレーキングには複数の機構がある。封入空気の圧縮はその一つであるが、膨潤などの説明もあり、全ての岩石について一つのモデルへ還元しない\cite{slaking1979,slaking2006}。
閉じ込められた空気の体積が減れば空気の圧力が上がる。しかし固体の孔壁は内部から押し広げられるため、骨格には引張破壊を促す応力が生じる。例えば単純化した薄い球殻では内外圧差\(\Delta p\)、半径\(r\)、厚さ\(t\)に対して膜の引張応力は\(\sigma\simeq\Delta p r/(2t)\)となる。この式は岩石の実形状への予測式ではなく、空気の圧縮と骨格の圧縮を同じにしないための説明用モデルである。

\section{不均質・異方性と弱面の向き}
\label{theory:engineering:anisotropy}
メランジュは細粒の基質にさまざまな岩塊を含み、層としての連続性に乏しい。硬い岩塊と弱い基質の分布、割れ目、水の経路を確認する\cite{jseg_engineering_faq}。結晶片岩の片理面では力学応答が向きで変わる。面が斜面へ抜けるか、抵抗する向きかという幾何も重要である。
せん断抵抗の説明用モデルとして\(\tau_f=c'+(\sigma_n-u)\tan\phi'\)を考えると、同じ全垂直応力\(\sigma_n\)でも間隙水圧\(u\)が増すと有効応力が減る。\(c'\)は有効粘着力、\(\phi'\)は有効摩擦角である。弱面の強度と岩片の強度は同じとは限らず、この単純モデルだけで複雑な斜面の安全率を決められない。

\section{付加体の地盤リスクと解答不一致}
\label{theory:engineering:accretion}
付加・変形によって複雑な構造を持つ地域では、地下水や弱面の連続性が問題になる。メランジュ地域の深層崩壊は実際に報告されている\cite{kyoto_deep2018,sabo_deep2014}。しかし深層崩壊の事例があることから表層崩壊が起こり得ないとは導けない。
蛇紋岩を含む地山の問題は、変質や破砕・粘土化の程度に依存する。粘土質蛇紋岩の膨張性地圧等の事例は、全ての健全な蛇紋岩が吸水膨張するという主張ではない\cite{jsce_swelling2010,serpentinite_mechanics2003}。吸水膨潤と掘削後の塑性流動も区別する。
平成27年の付加体問題では独立判断5と公式1が一致しない。原典5の海山を斜面等へ改変せず、原典1の表層崩壊を否定する限定根拠も未確定である。問題幹・全肢・独立記録・公式表画像の照合を保存し、Astra査読事項として残す。

\section{過去問で確認すること}
土の名称なら母材・生成・運搬、岩盤なら材料強度・弱面・水・応力を分ける。一般的な傾向を全試料に通じる規則にしない。地盤調査・各種試験・対策の測定量と適用条件を次の共通理論から確認する。
\section{登録問題と解説初稿}
\section{岩片強度と岩盤の代理指標}
\label{theory:engineering:strengthindices}
一軸圧縮強度は側方を拘束しない岩石試験片の最大圧縮応力である。岩盤では割れ目、弱面、水、寸法効果が加わり、試験片の値をそのまま現地全体へ使わない。RQDは所定長以上の健全コアの長さ割合に基づく岩盤状態の指標で、打診音は硬さ・風化・割れ目等の現地判別である。弾性波速度は弾性率と密度・接触状態を反映する。いずれも強度を一意に測る式ではなく、調査・試験を組み合わせて判断する\cite{jseg_engineering_faq,geebec_ps}。
圧密降伏応力は通常、土の圧縮履歴や降伏を表す量である。令和元年度再試験問27では岩盤強度との関連が最も薄いものとして選ぶ。ただし弱い泥岩等の圧縮降伏との関連まで否定しない。物理探査のデータを含む設問でも、中心は岩盤の強度評価であるため本章へ置く。
\section{三相と分母：何に対する割合か}\label{theory:engineering:phases}
土粒子の体積を\(V_s\)、水を\(V_w\)、空気を\(V_a\)とし、
全体積を\(V=V_s+V_w+V_a\)、間隙体積を\(V_v=V_w+V_a\)とする。
乾燥土粒子の質量を\(M_s\)、水質量を\(M_w\)とすれば
\[
 w=\frac{M_w}{M_s},\quad n=\frac{V_v}{V},\quad
 e=\frac{V_v}{V_s},\quad \theta=\frac{V_w}{V},\quad
 S_r=\frac{V_w}{V_v}.
\]
順に含水比、間隙率、間隙比、体積含水率、飽和度である。
百分率で表す場合は該当する比に100を掛ける。\(n=e/(1+e)\)、\(\theta=nS_r\)は
同じ三相の体積から導ける。含水比は水が乾燥土より重ければ1を超えるが、
体積の定義に基づく\(n,\theta,S_r\)は通常0から1である。\cite{eng52_soil_phases}

湿潤質量\(M_s+M_w\)で割る含水率と、乾燥質量で割る土質の含水比を取り違えない。
質量を体積で割れば密度の次元が残るので、百分率の候補を除外できる。
乾燥質量は測定条件に依存する。平成25年度の指定\(110\pm5\,^\circ\mathrm C\)を保持し、
別の試験規約の温度へ勝手に置き換えない。

\section{有効応力から液状化を説明する}\label{theory:engineering:liquefaction}
飽和土の骨格へ作用する応力を、説明用に\(\sigma'=\sigma-u\)と表す。
全応力\(\sigma\)と間隙水圧\(u\)を分ける。緩い砂質土が短時間の繰返しせん断を受け、
水が十分に排出できない場合、収縮しようとする骨格の応答が過剰間隙水圧へ現れる。
uが上がると有効応力が下がり、抵抗・剛性が低下して大きな変形が生じ得る。\cite{eng52_liquefaction}

高い地下水位、若く緩い砂質堆積物、よく淘汰された粒度等は代表的な条件である。
旧河道・埋立地・浜堤列間の低地という地形履歴を地下水・粒度・密度へつなぐ。
N値は標準貫入の抵抗の指標で、低いNは緩い砂の傾向を示すが、上載圧や細粒分・地震外力も考える。
原典のN20以上やN10以下を全土層の安全・危険を分ける一つの絶対線とはしない。
淘汰の悪い土も組成・密度によって挙動が異なる。\cite{eng52_liq_sorting}

試料を採取して動的な変形・液状化抵抗を調べ、PS検層で応答に関係する波速を得る。
試料の乱れと現場の密度を確認する。通常の平板載荷は静的な支持・沈下が主目的なので、
発生の評価の相対比較では直接関係が薄い。しかし液状化後の地盤反力を平板載荷で測る研究もある。
「最も関係が薄い」を「絶対に利用されない」へ変えない。\cite{eng52_plate_liq,geebec_ps}

\section{試験名を測定量と条件へ戻す}\label{theory:engineering:tests}
三軸試験は拘束圧と排水・圧密の条件を指定し、破壊と変形を測る。
一軸試験は側方拘束を与えずに圧縮する。ベーン試験は回転のせん断抵抗、
標準貫入やサウンディングは貫入抵抗を調べる。
前二者の破壊試験と、貫入から経験的に推定する強度・支持性を同じ直接測定とはしない。\cite{eng52_geotech_tests}
圧密試験は荷重の増加に対する時間依存の体積変化・圧縮性・排水を調べる。
状態が変われば強度も変わり得るが、主目的はせん断破壊強度の直接測定ではない。\cite{eng52_consolidation}

\section{Mohr--Coulombの切片・傾きと測定丸め}\label{theory:engineering:mohr}
原典が直線の包絡線を指定したときは
\[
 \tau_f=c+\sigma_n\tan\phi
\]
として切片cと傾き\(\tan\phi\)を分ける。
有効応力の強度定数を使う場合は
\(\tau_f=c'+(\sigma_n-u)\tan\phi'\)となるが、問題の記号を無断で有効応力へ変更しない。
摩擦角の値をそのまま傾きへ入れず、正接を使う。\cite{eng52_mohr_ref}

平成29年度の測定は共通単位\(10^4\,\mathrm{Pa}\)で
\((1,3.1),(2,3.7),(4,4.8),(8,7.1)\)。
端点の傾きは\(4/7\simeq0.5714\)で指定\(\tan30^\circ=0.58\)に近い。
各点の\(\tau-0.58\sigma\)は2.52、2.54、2.48、2.46で、切片2.5の候補と整合する。
丸めた測定点は厳密には一直線ではない。通常の最小二乗でも
傾き0.5696、切片2.5391となり同じ選択となる。数値単位はcにも掛かる。
残差をゼロにする架空の測定値へ修正しない。

\section{揚水と圧密：水が抜けて骨格へ荷重が移る}\label{theory:engineering:consolidation}
地下水揚水で被圧帯水層の水頭が下がると、隣接する粘土層の水圧にも影響する。
同じ全応力でuが低下すれば有効応力は増加する。透水性が低い粘土では、
排水に時間を要して収縮し、地表の沈下として現れる。\cite{eng52_consolidation,eng52_subsidence_ref}
砂の即時的な変形、粘土の圧密、地震による液状化後の体積変化を区別する。
水頭低下と沈下が同時同量でないことにも注意する。

杭等に支えられた構造物の相対的な浮き上がり、道路・埋設物の変形に加え、
低平地では排水・浸水リスクも変わる。深度ごとに帯水層の水頭を観測し、
複数層をつないだ混合水位だけで原因を判断しない。
不圧地下水位の低下だけを広域沈下の唯一の機構としない。

\section{黄鉄鉱の酸化と、溶解・析出の時間順}\label{theory:engineering:pyrite}
還元的な環境で保存された黄鉄鉱は、掘削後に水と酸素へさらされると酸化する。
例えば二価鉄を生成する段階を収支式で書くと
\[
 \mathrm{FeS_2}+\frac72\mathrm{O_2}+\mathrm{H_2O}
 \longrightarrow \mathrm{Fe^{2+}}+2\mathrm{SO_4^{2-}}+2\mathrm{H^+}.
\]
さらに鉄の酸化・加水分解、炭酸塩による中和等が加わるので、実際のpHは材料と水の条件に依存する。
酸化そのものをアルカリ性排水の起源とはしない。\cite{eng52_weathering}

酸性水が膠結物等を溶かして岩石を弱める過程と、溶けた成分が石膏等の硫酸塩で再析出する過程を分ける。
蒸発・過飽和・水の経路によって析出し、結晶成長圧が盤膨れへつながる例がある。
青黒い未酸化岩の性質を、掘削後も酸素に触れないまま硫黄の酸化が進むという説明へ変えない。
酸化帯の充填、下の溶解帯の弱化などがあるため、酸化の存在だけで強化又は抑制を一方向へ決めない。\cite{eng52_weathering}

\section{鉄の色・還元溶出と交換サイト}\label{theory:engineering:soilredox}
酸化的な鉄の酸化物等は赤褐色・黄褐色、還元的なグライ化は青灰色等という代表的な対応がある。
土壌の色は腐植等にも依存するので一つの鉄化合物だけから全部を決めない。
湛水して酸素が消費されるとFe・Mnが還元され、二価の形で溶出しやすくなる場合がある。\cite{eng52_soil_colour}
水圧による力学的な不安定化と、酸化還元による化学的な物性変化を分ける。

粘土・腐植の交換サイトではカチオンが電荷を保ちながら入れ替わる。
同じ固定サイト系での保存条件と、pH等に伴う材料の交換容量の変化を区別する。
土壌有機炭素は重要だが、炭酸塩に富む土壌を含む全炭素を有機物だけへ限定しない。

\section{斜面：弱面と地下水を幾何へ重ねる}\label{theory:engineering:landslide}
重力性の移動を物質・移動様式・速度の履歴で記述する。
層理面・変質凝灰岩等の弱層が斜面へ抜ける流れ盤ではすべりへ注意する。
背斜・ドーム・キャップロックという名前だけで全斜面の安定を決めず、割れ目・水・地形を重ねる。\cite{eng52_slip_weaklayer}
すべり面の抵抗を\(c'+(\sigma_n-u)\tan\phi'\)で考えると、
水圧の上昇で抵抗が下がり得ることを説明できる。

滑落崖、亀裂、隆起・陥没は移動を示す観察対象である。長期緩慢・反復する地すべりも、
急激な滑動もある。「地すべり＝全例低速、崩壊＝全例脆性」の二分法は使わない。
原典の砂質/粘性の逆転は典型比較の中で説明し、全岩種を排他的に分類する根拠にはしない。\cite{eng52_landslide_ref}
浸透を減らす被覆・地下水排水・押え盛土等は機構に応じた対策候補であり、
現地の安全と施工条件を確認せず実行する一般指示ではない。

\section{トンネル：切羽・緩み・応力集中}\label{theory:engineering:tunnel}
掘削は空洞周囲の応力を変え、不連続面に沿う落下・緩み、水による切羽の弱化等を生じ得る。
切羽の天端・鏡面を観察し、地質構造と湧水を支保・設計へ反映する。\cite{eng52_tunnel_ref}
山はねは代表的には高応力・大きな土かぶり・硬く節理の少ない岩盤で問題になり、
小土かぶり・多節理を典型条件とする原典の候補は逆である。\cite{eng52_rockburst}
斜面地形・傾斜地層等の非対称な偏圧、浅部の地表沈下、地下水低下による沈下を分ける。
深さだけで沈下量を一つの単調則として決めない。

\section{骨材：岩種は品質試験の入口}\label{theory:engineering:aggregate}
岩石名から鉱物組成・強度・耐久性・吸水・反応性を考えるが、名称だけで品質を保証しない。
蛇紋岩骨材のポップアウト研究はある一方、骨材用採石の記録もある。
公式の一般的な不適性の比較と、全試料の使用不可能という全称を区別する。\cite{eng52_serpent_aggregate,eng52_popout}
チャートのシリカの反応性、砂岩の結合、片麻岩の方向性等も考える。
一つの健全試料の結果を全変質・風化状態へ広げない。

\section{浅部地中熱：伝導と移流を分離する}\label{theory:engineering:groundheat}
浅部の安定した温度と外気との差を利用して冷暖房する地中熱を、深部の高温地熱発電と区別する。
閉じた地中配管で媒体を循環するクローズドループと、地下水を利用するオープンループを分ける。
地盤・帯水層を熱の貯蔵へ使う場合は蓄熱の履歴と周辺環境への影響も考える。\cite{eng52_groundheat_ref}

熱伝導は温度勾配による輸送、地下水の移流は水が熱を運ぶ過程である。
熱応答試験を伝導モデルで整理すると移流の効果が見掛けの有効熱伝導率へ現れ、
材料自体の伝導率より大きく推定される例がある。\cite{eng52_heat_advection}
原典の一般的な大幅減少という向きを選ぶ誤りとして区別する。
その寄与は流速・方向・幾何・試験時間に依存し、全試験の一定倍率ではない。

\section{付加体の構成岩・形成過程}

\subsection{平成27年度 問34}
\label{q:h27-reg-applied-34}

\noindent\textbf{原典：} sources/H27/17\_applied\_science.pdf、PDF物理13頁。

\textbf{公式正答：} 1（single）。

\par\textbf{独立解答と公式正答に相違あり：} 独立解答 5。初回独立解答5、公式1。公式表107頁を実画像で再確認。1の表層崩壊を否定する排他的根拠は未確定。5の大陸側斜面と海山の関係も原典のまま保持。メランジュの深層崩壊事例を表層崩壊の不可能性へ読み替えない。Astraで原典と地質機構を重点査読。

\paragraph{問題の要旨} 付加体の構造・弱面と斜面安定性に関する記述を判定する。

\Needspace{4\baselineskip}\paragraph{確認した条件・記述（独自要約）}\mbox{}\par

\noindent\textbf{方向：} 最も不適切なもの。\par

\noindent\textbf{対象：} 一般的な岩石分類・形成過程。\par

\paragraph{問われる核心} 付加体の不均質・異方性・斜面崩壊と未解決の肢解釈

\paragraph{必要知識} 付加体の不均質・異方性・斜面崩壊と未解決の肢解釈。母材・運搬堆積・有機物起源、風化、弱面、水の作用、岩盤と土の物性。

\paragraph{考え方・正答を選ぶ理由（査読待ち初稿）} 公式正答は1であり、公式正答表107頁のIII–34の1を画像でも確認した。初回独立解答は5である。原典1の「表層崩壊」と5の「大陸側斜面」「海山」の限定を含め、公式1を支持する排他的理由は現在の一次資料照合では確定していない。メランジュに深層崩壊事例があることだけから1の表層崩壊を誤りとは導かない。相違を未解決として保持し、公式を独立の推論結果と偽らない。

\Needspace{4\baselineskip}\paragraph{選択肢の検討（査読待ち初稿）}\mbox{}\par

\noindent\textbf{1：公式では不適切／根拠未確定} 砂岩の岩塊を剥離性の泥質基質が囲むメランジュは表層崩壊等を生じやすい。。公式はこの肢を不適切とする。しかしメランジュの硬軟混在と割れ目が切土等で問題となることは学会資料で確認できる。深層崩壊との関連もあるが、深層の事例だけで表層崩壊を不可能とは言えない。誤りとする限定根拠は未確定。\par

\noindent\textbf{2：概ね適切} 片理による結晶片岩の力学的異方性は地盤安定性に影響する。。面に沿うせん断抵抗などが方向で変わり、弱面と斜面の位置関係を評価する必要がある。等方な単一強度では表せない。\par

\noindent\textbf{3：条件付きで適切} 蛇紋岩には膨張性地山や光沢ある鱗片状剥離面を持ち、地すべりを起こしやすいものがある。。蛇紋岩地域の地すべりや粘土化した蛇紋岩の膨張性地山は一次資料に事例がある。健全な全蛇紋岩が膨潤するとの主張にはしない。\par

\noindent\textbf{4：概ね適切} 海溝の堆積物は圧縮で形成された衝上断層に沿って大陸側へ付加される。。付加は海洋プレート上の堆積物等が剥ぎ取られ、陸側に取り込まれる過程で、衝上断層が発達する。侵食型の沈み込み域もある。\par

\noindent\textbf{5：初回独立判断で疑義／未確定} 海溝の大陸側斜面で海底地すべりが起こり、海山が断裂して海溝へ崩壊する。。海底斜面崩壊自体は起こり得る。一方、海山は一般に沈み込む海洋プレート上の地形で、この一文の大陸側斜面との関係を疑って初回独立解答5とした。沈み込み中の海山破壊も起こり得るため、単に海山という語だけで絶対的な誤りとは確定しない。\par

\paragraph{再利用する考え方} 付加体の不均質・異方性・斜面崩壊と未解決の肢解釈を成因と観察される性質に結び付け、例示と全称命題、主要鉱物と副成分、公式解答と未解決の解釈を区別する。

\paragraph{共通理論} \hyperref[theory:engineering:accretion]{第\ref*{theory:engineering:accretion}節}, \hyperref[theory:engineering:anisotropy]{第\ref*{theory:engineering:anisotropy}節}

\par 参考：\cite{jseg_engineering_faq,kyoto_deep2018,sabo_deep2014,jsce_swelling2010,aist_accretion2026}。

\par 難易度 3（校準初稿）、優先度未評価。

\section{骨材の岩種・変質と品質}

\subsection{平成24年度 問28}
\label{q:h24-reg-applied-28}

\noindent\textbf{原典：} sources/H24/17\_applied\_science.pdf、PDF物理8頁。

\textbf{公式正答：} 2（single）。

\par\textbf{検討未確定：} 蛇紋岩のポップアウト研究は土研報告の参考文献で確認し、研究原報告の全文は未取得。各岩種の試験当時の総合順位と全試料の適性は未確定。

\paragraph{問題の要旨} 岩石名だけでコンクリート骨材の適性を固定しない。

\Needspace{4\baselineskip}\paragraph{確認した条件・記述（独自要約）}\mbox{}\par

\noindent\textbf{原典照合：} 22原典画像で25問の全五肢・数値・単位・否定を確認。\par

\noindent\textbf{対象：} 岩石名だけでコンクリート骨材の適性を固定しない。\par

\paragraph{問われる核心} 岩石名だけでコンクリート骨材の適性を固定しない。

\paragraph{必要知識} 選択2は一般的な骨材の相対比較として保持する。蛇紋岩によるポップアウトの報告や品質問題はあるが、厚労省の一次資料には骨材用採石も記載される。したがって全蛇紋岩の使用が絶対不可能という解説にはしない。

\paragraph{考え方・正答を選ぶ理由（査読待ち初稿）} 選択2は一般的な骨材の相対比較として保持する。蛇紋岩によるポップアウトの報告や品質問題はあるが、厚労省の一次資料には骨材用採石も記載される。したがって全蛇紋岩の使用が絶対不可能という解説にはしない。

\Needspace{4\baselineskip}\paragraph{選択肢の検討（査読待ち初稿）}\mbox{}\par

\noindent\textbf{1：条件・全称を留保} チャート。硬い珪質岩だが、微細なシリカを含む骨材ではアルカリシリカ反応等に注意。全チャートが無条件に安全とはしない。\par

\noindent\textbf{2：不適切} 蛇紋岩。公式・独立とも2。蛇紋岩骨材に起因するポップアウトの研究例がある。風化・変質・鉱物組成により強度・耐久性の適性が変わる。\par

\noindent\textbf{3：代表的条件で適切} 硬質砂岩。硬く耐久性を持つ適切な試料は骨材候補。砂岩という名称だけで粒子結合や吸水率等を省略しない。\par

\noindent\textbf{4：代表的条件で適切} 斑れい岩。健全で緻密な試料は骨材候補。鉱物組成・変質・有害成分等の試験も必要。\par

\noindent\textbf{5：条件・全称を留保} 片麻岩。片理や異方性を含めて品質評価する。すべての方向・試料で同じ破砕形状・強度とはしない。\par

\paragraph{再利用する考え方} 直接測る量と代理指標を区別し、分母・有効応力・水の経路・材料の条件から判断する。

\paragraph{共通理論} \hyperref[theory:engineering:aggregate]{第\ref*{theory:engineering:aggregate}節}

\par 参考：\cite{eng52_serpent_aggregate,eng52_popout}。

\par 難易度 3（校準初稿）、優先度未評価。

\section{浅部地中熱の方式・移流}

\subsection{令和6年度 問18}
\label{q:r06-reg-applied-18}

\noindent\textbf{原典：} sources/R06/17\_applied\_science.pdf、PDF物理6頁。

\textbf{公式正答：} 5（single）。

\par\textbf{検討未確定：} 浅度200m・低温・騒音は代表的な説明であり全機器・地点の厳密境界ではない。移流の寄与は流速・方向・試験期間等に依存し全例の同じ倍率ではない。

\paragraph{問題の要旨} 浅部地中熱の循環方式と移流による熱輸送を区別する。

\Needspace{4\baselineskip}\paragraph{確認した条件・記述（独自要約）}\mbox{}\par

\noindent\textbf{原典照合：} 22原典画像で25問の全五肢・数値・単位・否定を確認。\par

\noindent\textbf{対象：} 浅部地中熱の循環方式と移流による熱輸送を区別する。\par

\paragraph{問われる核心} 浅部地中熱の循環方式と移流による熱輸送を区別する。

\paragraph{必要知識} 誤り5は地下水の移流による熱輸送の向き。実際の熱伝導率は材料物性、流動を含む見掛け値はモデルと観測条件を反映する。熱応答試験で推定する有効値を岩石自体の不変な熱伝導率と同一視しない。

\paragraph{考え方・正答を選ぶ理由（査読待ち初稿）} 誤り5は地下水の移流による熱輸送の向き。実際の熱伝導率は材料物性、流動を含む見掛け値はモデルと観測条件を反映する。熱応答試験で推定する有効値を岩石自体の不変な熱伝導率と同一視しない。

\Needspace{4\baselineskip}\paragraph{選択肢の検討（査読待ち初稿）}\mbox{}\par

\noindent\textbf{1：代表的条件で適切} 概ね200mより浅い低温地中熱。通常の地中熱利用の目安。深部の高温地熱発電と区別するが、200mを物理的な熱源の断絶境界とはしない。\par

\noindent\textbf{2：代表的条件で適切} 外部騒音が一般に小さい。空気熱交換部等の条件を考える。機器の音が絶対にゼロではない。\par

\noindent\textbf{3：代表的条件で適切} クローズド・オープンループ。地中の閉じた配管で媒体を循環する方式と、地下水を利用する方式を分ける。\par

\noindent\textbf{4：代表的条件で適切} 地盤を蓄熱体にする地下蓄熱。季節等の熱を地盤又は帯水層へ蓄えて利用する。すべての採熱装置が蓄熱運用という意味ではない。\par

\noindent\textbf{5：不適切} 移流で見掛け熱伝導率が大きく減少。地下水移流は熱を運び、熱応答試験の伝導モデルで見掛けの有効熱伝導率が増える場合がある。一般的な減少効果という向きが不適切。\par

\paragraph{再利用する考え方} 直接測る量と代理指標を区別し、分母・有効応力・水の経路・材料の条件から判断する。

\paragraph{共通理論} \hyperref[theory:engineering:groundheat]{第\ref*{theory:engineering:groundheat}節}

\par 参考：\cite{eng52_groundheat_ref,eng52_heat_advection}。

\par 難易度 3（校準初稿）、優先度未評価。

\section{弱面・水圧・地すべり滑動}

\subsection{平成26年度 問33}
\label{q:h26-reg-applied-33}

\noindent\textbf{原典：} sources/H26/17\_applied\_science.pdf、PDF物理13頁。

\textbf{公式正答：} 2（single）。

\par\textbf{検討未確定：} 列挙された全構造の発生頻度や全応急工法の適性は追加照合対象。岩種・構造名だけで安定性又は対策の安全性を保証しない。

\paragraph{問題の要旨} 地すべりの弱面・水・滑動と対策を対応させる。

\Needspace{4\baselineskip}\paragraph{確認した条件・記述（独自要約）}\mbox{}\par

\noindent\textbf{原典照合：} 22原典画像で25問の全五肢・数値・単位・否定を確認。\par

\noindent\textbf{対象：} 地すべりの弱面・水・滑動と対策を対応させる。\par

\paragraph{問われる核心} 地すべりの弱面・水・滑動と対策を対応させる。

\paragraph{必要知識} 公式2は各回で誤りの内容が変わる。H26は構造との関係を否定、R01再試験は泥岩すべりの長期緩慢・反復を否定した点が不適切。地質構造、間隙水圧、滑落崖、速度の履歴を組み合わせて読む。

\paragraph{考え方・正答を選ぶ理由（査読待ち初稿）} 公式2は各回で誤りの内容が変わる。H26は構造との関係を否定、R01再試験は泥岩すべりの長期緩慢・反復を否定した点が不適切。地質構造、間隙水圧、滑落崖、速度の履歴を組み合わせて読む。

\Needspace{4\baselineskip}\paragraph{選択肢の検討（査読待ち初稿）}\mbox{}\par

\noindent\textbf{1：代表的条件で適切} 重力によるマスムーブメント。地すべり・崩壊・土石流・落石は斜面下方への重力性移動に関係する。移動様式は区別する。\par

\noindent\textbf{2：不適切} 流れ盤・背斜・ドーム・キャップで発生しにくい。斜面へ抜ける層理等の弱面、構造による割れ目、水の経路は発生条件になり得る。列挙した構造をまとめて発生しにくいとはしない。\par

\noindent\textbf{3：代表的条件で適切} 最上部の滑落崖・馬蹄形。移動体上端の崖を滑落崖と呼ぶ。馬蹄形はよくある形で、全崖の唯一形ではない。\par

\noindent\textbf{4：代表的条件で適切} 泥岩で緩慢・長期、雨季・融雪で増加。地下水変動や弱いすべり面に対応した反復・継続の例がある。全地すべりの速度を一定にしない。\par

\noindent\textbf{5：条件・全称を留保} 浸透防止・横孔排水・押え盛土。水圧を下げたり滑動を抑えたりする対策。応急対策の施工位置・安全性とすべりの機構の確認が必要で、すべてを無条件に現場実施する指示ではない。\par

\paragraph{再利用する考え方} 直接測る量と代理指標を区別し、分母・有効応力・水の経路・材料の条件から判断する。

\paragraph{共通理論} \hyperref[theory:engineering:landslide]{第\ref*{theory:engineering:landslide}節}

\par 参考：\cite{eng52_slip_weaklayer,eng52_landslide_ref}。

\paragraph{関連問題} \hyperref[q:r01-ret-applied-26]{R01-RET-APPLIED-26}

\par 難易度 3（校準初稿）、優先度未評価。

\subsection{令和元年度再試験 問26}
\label{q:r01-ret-applied-26}

\noindent\textbf{原典：} sources/R01\_retest/17\_applied\_science.pdf、PDF物理9頁。

\textbf{公式正答：} 2（single）。

\par\textbf{検討未確定：} 列挙された全構造の発生頻度や全応急工法の適性は追加照合対象。岩種・構造名だけで安定性又は対策の安全性を保証しない。

\paragraph{問題の要旨} 地すべりの弱面・水・滑動と対策を対応させる。

\Needspace{4\baselineskip}\paragraph{確認した条件・記述（独自要約）}\mbox{}\par

\noindent\textbf{原典照合：} 22原典画像で25問の全五肢・数値・単位・否定を確認。\par

\noindent\textbf{対象：} 地すべりの弱面・水・滑動と対策を対応させる。\par

\paragraph{問われる核心} 地すべりの弱面・水・滑動と対策を対応させる。

\paragraph{必要知識} 公式2は各回で誤りの内容が変わる。H26は構造との関係を否定、R01再試験は泥岩すべりの長期緩慢・反復を否定した点が不適切。地質構造、間隙水圧、滑落崖、速度の履歴を組み合わせて読む。

\paragraph{考え方・正答を選ぶ理由（査読待ち初稿）} 公式2は各回で誤りの内容が変わる。H26は構造との関係を否定、R01再試験は泥岩すべりの長期緩慢・反復を否定した点が不適切。地質構造、間隙水圧、滑落崖、速度の履歴を組み合わせて読む。

\Needspace{4\baselineskip}\paragraph{選択肢の検討（査読待ち初稿）}\mbox{}\par

\noindent\textbf{1：代表的条件で適切} 重力によるマスムーブメント。地すべり・崩壊・土石流・落石は斜面下方への重力性移動に関係する。移動様式は区別する。\par

\noindent\textbf{2：不適切} 泥岩で長期反復が少なく短期急激。典型的な泥岩地域の地すべりには長期緩慢・反復的な滑動があり、その特徴を逆にした記述。急激な移動の例まで否定するものではない。\par

\noindent\textbf{3：条件・全称を留保} 流れ盤等で発生しやすい。原典の代表傾向として読む。弱面が斜面へ抜ける角度、強度、水、地形との組合せが必要で、背斜名だけから全地点を決めない。\par

\noindent\textbf{4：代表的条件で適切} 最上部の滑落崖・馬蹄形。移動体上端の崖を滑落崖と呼ぶ。馬蹄形はよくある形で、全崖の唯一形ではない。\par

\noindent\textbf{5：条件・全称を留保} 浸透防止・横孔排水・押え盛土。水圧を下げたり滑動を抑えたりする対策。応急対策の施工位置・安全性とすべりの機構の確認が必要で、すべてを無条件に現場実施する指示ではない。\par

\paragraph{再利用する考え方} 直接測る量と代理指標を区別し、分母・有効応力・水の経路・材料の条件から判断する。

\paragraph{共通理論} \hyperref[theory:engineering:landslide]{第\ref*{theory:engineering:landslide}節}

\par 参考：\cite{eng52_slip_weaklayer,eng52_landslide_ref}。

\paragraph{関連問題} \hyperref[q:h26-reg-applied-33]{H26-REG-APPLIED-33}

\par 難易度 3（校準初稿）、優先度未評価。

\section{飽和・粒度・密度の液状化条件}

\subsection{平成24年度 問31}
\label{q:h24-reg-applied-31}

\noindent\textbf{原典：} sources/H24/17\_applied\_science.pdf、PDF物理9頁。

\textbf{公式正答：} 5（single）。

\par\textbf{検討未確定：} N20を全地震動・全土層の不液状化保証へ使わない。R03の淘汰の悪さも唯一の排他的条件ではなく、原典は代表的な相対比較として保持。

\paragraph{問題の要旨} 液状化しやすい飽和・粒度・密度の条件を比べる。

\Needspace{4\baselineskip}\paragraph{確認した条件・記述（独自要約）}\mbox{}\par

\noindent\textbf{原典照合：} 22原典画像で25問の全五肢・数値・単位・否定を確認。\par

\noindent\textbf{対象：} 液状化しやすい飽和・粒度・密度の条件を比べる。\par

\paragraph{問われる核心} 液状化しやすい飽和・粒度・密度の条件を比べる。

\paragraph{必要知識} 緩い粒子骨格が飽和し、繰返しせん断で水圧が上がる条件を考える。公式5はこの代表比較で最も不適切な候補。土地名・N値・淘汰の一条件だけで判定を完結せず、密度・粒度・地下水と地震外力を組み合わせる。

\paragraph{考え方・正答を選ぶ理由（査読待ち初稿）} 緩い粒子骨格が飽和し、繰返しせん断で水圧が上がる条件を考える。公式5はこの代表比較で最も不適切な候補。土地名・N値・淘汰の一条件だけで判定を完結せず、密度・粒度・地下水と地震外力を組み合わせる。

\Needspace{4\baselineskip}\paragraph{選択肢の検討（査読待ち初稿）}\mbox{}\par

\noindent\textbf{1：代表的条件で適切} 埋立地・旧河道。若く緩い砂質堆積物や高い地下水位がある場所ではリスクが高まり得る。名称だけで全地点を判定しない。\par

\noindent\textbf{2：代表的条件で適切} 砂丘・浜堤列の間の低地。地下水が浅く飽和した砂質地盤が分布し得る。高い砂丘そのものと列間低地を区別する。\par

\noindent\textbf{3：代表的条件で適切} 高い地下水位。浅い飽和砂層が揺れを受ける条件を作る。水位だけで粒度・密度・揺れを省略しない。\par

\noindent\textbf{4：代表的条件で適切} 粒度が一定の堆積物。よく淘汰された緩い飽和砂は代表的に液状化しやすい。粒度だけで全密度の判定をしない。\par

\noindent\textbf{5：条件・全称を留保} N値20以上。密度・締まりが相対的に高い傾向で、緩い砂の典型条件との関係は小さい。地震動・有効上載圧・粒度等の補正があり、20を絶対的な安全線とはしない。\par

\paragraph{再利用する考え方} 直接測る量と代理指標を区別し、分母・有効応力・水の経路・材料の条件から判断する。

\paragraph{共通理論} \hyperref[theory:engineering:liquefaction]{第\ref*{theory:engineering:liquefaction}節}

\par 参考：\cite{eng52_liquefaction,eng52_liq_sorting}。

\paragraph{関連問題} \hyperref[q:h27-reg-applied-35]{H27-REG-APPLIED-35}

\par 難易度 3（校準初稿）、優先度未評価。

\subsection{平成27年度 問35}
\label{q:h27-reg-applied-35}

\noindent\textbf{原典：} sources/H27/17\_applied\_science.pdf、PDF物理14頁。

\textbf{公式正答：} 5（single）。

\par\textbf{検討未確定：} N20を全地震動・全土層の不液状化保証へ使わない。R03の淘汰の悪さも唯一の排他的条件ではなく、原典は代表的な相対比較として保持。

\paragraph{問題の要旨} 液状化しやすい飽和・粒度・密度の条件を比べる。

\Needspace{4\baselineskip}\paragraph{確認した条件・記述（独自要約）}\mbox{}\par

\noindent\textbf{原典照合：} 22原典画像で25問の全五肢・数値・単位・否定を確認。\par

\noindent\textbf{対象：} 液状化しやすい飽和・粒度・密度の条件を比べる。\par

\paragraph{問われる核心} 液状化しやすい飽和・粒度・密度の条件を比べる。

\paragraph{必要知識} 緩い粒子骨格が飽和し、繰返しせん断で水圧が上がる条件を考える。公式5はこの代表比較で最も不適切な候補。土地名・N値・淘汰の一条件だけで判定を完結せず、密度・粒度・地下水と地震外力を組み合わせる。

\paragraph{考え方・正答を選ぶ理由（査読待ち初稿）} 緩い粒子骨格が飽和し、繰返しせん断で水圧が上がる条件を考える。公式5はこの代表比較で最も不適切な候補。土地名・N値・淘汰の一条件だけで判定を完結せず、密度・粒度・地下水と地震外力を組み合わせる。

\Needspace{4\baselineskip}\paragraph{選択肢の検討（査読待ち初稿）}\mbox{}\par

\noindent\textbf{1：代表的条件で適切} 埋立地・旧河道。若く緩い砂質堆積物や高い地下水位がある場所ではリスクが高まり得る。名称だけで全地点を判定しない。\par

\noindent\textbf{2：代表的条件で適切} 砂丘・浜堤列の間の低地。地下水が浅く飽和した砂質地盤が分布し得る。高い砂丘そのものと列間低地を区別する。\par

\noindent\textbf{3：代表的条件で適切} 高い地下水位。浅い飽和砂層が揺れを受ける条件を作る。水位だけで粒度・密度・揺れを省略しない。\par

\noindent\textbf{4：代表的条件で適切} 粒度が一定の堆積物。よく淘汰された緩い飽和砂は代表的に液状化しやすい。粒度だけで全密度の判定をしない。\par

\noindent\textbf{5：条件・全称を留保} N値20以上。密度・締まりが相対的に高い傾向で、緩い砂の典型条件との関係は小さい。地震動・有効上載圧・粒度等の補正があり、20を絶対的な安全線とはしない。\par

\paragraph{再利用する考え方} 直接測る量と代理指標を区別し、分母・有効応力・水の経路・材料の条件から判断する。

\paragraph{共通理論} \hyperref[theory:engineering:liquefaction]{第\ref*{theory:engineering:liquefaction}節}

\par 参考：\cite{eng52_liquefaction,eng52_liq_sorting}。

\paragraph{関連問題} \hyperref[q:h24-reg-applied-31]{H24-REG-APPLIED-31}, \hyperref[q:r03-reg-applied-29]{R03-REG-APPLIED-29}

\par 難易度 3（校準初稿）、優先度未評価。

\subsection{令和3年度 問29}
\label{q:r03-reg-applied-29}

\noindent\textbf{原典：} sources/R03/17\_applied\_science.pdf、PDF物理10頁。

\textbf{公式正答：} 2（single）。

\par\textbf{検討未確定：} N20を全地震動・全土層の不液状化保証へ使わない。R03の淘汰の悪さも唯一の排他的条件ではなく、原典は代表的な相対比較として保持。

\paragraph{問題の要旨} 液状化しやすい飽和・粒度・密度の条件を比べる。

\Needspace{4\baselineskip}\paragraph{確認した条件・記述（独自要約）}\mbox{}\par

\noindent\textbf{原典照合：} 22原典画像で25問の全五肢・数値・単位・否定を確認。\par

\noindent\textbf{対象：} 液状化しやすい飽和・粒度・密度の条件を比べる。\par

\paragraph{問われる核心} 液状化しやすい飽和・粒度・密度の条件を比べる。

\paragraph{必要知識} 緩い粒子骨格が飽和し、繰返しせん断で水圧が上がる条件を考える。公式2はこの代表比較で最も不適切な候補。土地名・N値・淘汰の一条件だけで判定を完結せず、密度・粒度・地下水と地震外力を組み合わせる。

\paragraph{考え方・正答を選ぶ理由（査読待ち初稿）} 緩い粒子骨格が飽和し、繰返しせん断で水圧が上がる条件を考える。公式2はこの代表比較で最も不適切な候補。土地名・N値・淘汰の一条件だけで判定を完結せず、密度・粒度・地下水と地震外力を組み合わせる。

\Needspace{4\baselineskip}\paragraph{選択肢の検討（査読待ち初稿）}\mbox{}\par

\noindent\textbf{1：代表的条件で適切} 高い地下水位。浅い飽和砂層が揺れを受ける条件を作る。水位だけで粒度・密度・揺れを省略しない。\par

\noindent\textbf{2：条件・全称を留保} 淘汰の悪い堆積物。均一粒径の緩い飽和砂と比べれば典型条件から離れる。ただし細粒分の種類・密度・組成によって液状化し得るため、淘汰だけで全例の発生を否定しない。\par

\noindent\textbf{3：代表的条件で適切} N値10以下。低い貫入抵抗は緩い砂の代表的な指標。粘性土の低N値と飽和砂の液状化機構を同一視しない。\par

\noindent\textbf{4：代表的条件で適切} 埋立地・旧河道。若く緩い砂質堆積物や高い地下水位がある場所ではリスクが高まり得る。名称だけで全地点を判定しない。\par

\noindent\textbf{5：代表的条件で適切} 砂丘・浜堤列の間の低地。地下水が浅く飽和した砂質地盤が分布し得る。高い砂丘そのものと列間低地を区別する。\par

\paragraph{再利用する考え方} 直接測る量と代理指標を区別し、分母・有効応力・水の経路・材料の条件から判断する。

\paragraph{共通理論} \hyperref[theory:engineering:liquefaction]{第\ref*{theory:engineering:liquefaction}節}

\par 参考：\cite{eng52_liquefaction,eng52_liq_sorting}。

\paragraph{関連問題} \hyperref[q:h27-reg-applied-35]{H27-REG-APPLIED-35}

\par 難易度 3（校準初稿）、優先度未評価。

\section{液状化の動的評価と静的載荷}

\subsection{平成23年度 問33}
\label{q:h23-reg-applied-33}

\noindent\textbf{原典：} sources/H23/17\_applied\_science.pdf、PDF物理11頁。

\textbf{公式正答：} 5（single）。

\par\textbf{検討未確定：} 公式選択は維持するが平板載荷と液状化を絶対的に無関係とする定義ではない。土研・研究論文には液状化後の反力を載荷で検討する例がある。

\paragraph{問題の要旨} 液状化の動的評価と静的な平板載荷を相対比較する。

\Needspace{4\baselineskip}\paragraph{確認した条件・記述（独自要約）}\mbox{}\par

\noindent\textbf{原典照合：} 22原典画像で25問の全五肢・数値・単位・否定を確認。\par

\noindent\textbf{対象：} 液状化の動的評価と静的な平板載荷を相対比較する。\par

\paragraph{問われる核心} 液状化の動的評価と静的な平板載荷を相対比較する。

\paragraph{必要知識} 選択肢5は通常の静的な平板載荷試験。地震時の繰返しせん断、密度を保った試料、地盤応答を調べる動的試験・PS検層と主目的が違う。液状化後の支持力等の研究では載荷試験も使われるため、原典の相対比較として答える。

\paragraph{考え方・正答を選ぶ理由（査読待ち初稿）} 選択肢5は通常の静的な平板載荷試験。地震時の繰返しせん断、密度を保った試料、地盤応答を調べる動的試験・PS検層と主目的が違う。液状化後の支持力等の研究では載荷試験も使われるため、原典の相対比較として答える。

\Needspace{4\baselineskip}\paragraph{選択肢の検討（査読待ち初稿）}\mbox{}\par

\noindent\textbf{1：代表的条件で適切} 乱さない試料採取。密度・粒子骨格等の状態をなるべく保持し、動的試験や液状化抵抗の評価へつなぐ。採取だけで発生の有無を直接決めるものではない。\par

\noindent\textbf{2：代表的条件で適切} 動的変形特性試験。繰返し載荷に対する剛性・減衰等を測り、地震応答を検討する。特性試験と液状化強度試験の目的は区別する。\par

\noindent\textbf{3：代表的条件で適切} PS検層。P・S波速度から地盤の動的性質・応答を検討する。波速だけで液状化を一意に判定しない。\par

\noindent\textbf{4：代表的条件で適切} 繰返しせん断力。排水が追いつかない飽和土で過剰間隙水圧が上がり、有効応力とせん断抵抗が低下する機構に関係する。\par

\noindent\textbf{5：条件・全称を留保} 平板載荷試験。通常は静的な支持力・荷重沈下関係を調べるため、候補中では液状化発生の評価との直接関係が最も薄い。ただし液状化後の地盤反力を載荷で調べた実験もあり、全く無関係とはしない。\par

\paragraph{再利用する考え方} 直接測る量と代理指標を区別し、分母・有効応力・水の経路・材料の条件から判断する。

\paragraph{共通理論} \hyperref[theory:engineering:liquefaction]{第\ref*{theory:engineering:liquefaction}節}

\par 参考：\cite{eng52_liquefaction,eng52_plate_liq,geebec_ps}。

\paragraph{関連問題} \hyperref[q:h28-reg-applied-33]{H28-REG-APPLIED-33}

\par 難易度 3（校準初稿）、優先度未評価。

\subsection{平成28年度 問33}
\label{q:h28-reg-applied-33}

\noindent\textbf{原典：} sources/H28/17\_applied\_science.pdf、PDF物理13頁。

\textbf{公式正答：} 1（single）。

\par\textbf{検討未確定：} 公式選択は維持するが平板載荷と液状化を絶対的に無関係とする定義ではない。土研・研究論文には液状化後の反力を載荷で検討する例がある。

\paragraph{問題の要旨} 液状化の動的評価と静的な平板載荷を相対比較する。

\Needspace{4\baselineskip}\paragraph{確認した条件・記述（独自要約）}\mbox{}\par

\noindent\textbf{原典照合：} 22原典画像で25問の全五肢・数値・単位・否定を確認。\par

\noindent\textbf{対象：} 液状化の動的評価と静的な平板載荷を相対比較する。\par

\paragraph{問われる核心} 液状化の動的評価と静的な平板載荷を相対比較する。

\paragraph{必要知識} 選択肢1は通常の静的な平板載荷試験。地震時の繰返しせん断、密度を保った試料、地盤応答を調べる動的試験・PS検層と主目的が違う。液状化後の支持力等の研究では載荷試験も使われるため、原典の相対比較として答える。

\paragraph{考え方・正答を選ぶ理由（査読待ち初稿）} 選択肢1は通常の静的な平板載荷試験。地震時の繰返しせん断、密度を保った試料、地盤応答を調べる動的試験・PS検層と主目的が違う。液状化後の支持力等の研究では載荷試験も使われるため、原典の相対比較として答える。

\Needspace{4\baselineskip}\paragraph{選択肢の検討（査読待ち初稿）}\mbox{}\par

\noindent\textbf{1：条件・全称を留保} 平板載荷試験。通常は静的な支持力・荷重沈下関係を調べるため、候補中では液状化発生の評価との直接関係が最も薄い。ただし液状化後の地盤反力を載荷で調べた実験もあり、全く無関係とはしない。\par

\noindent\textbf{2：代表的条件で適切} 動的変形特性試験。繰返し載荷に対する剛性・減衰等を測り、地震応答を検討する。特性試験と液状化強度試験の目的は区別する。\par

\noindent\textbf{3：代表的条件で適切} PS検層。P・S波速度から地盤の動的性質・応答を検討する。波速だけで液状化を一意に判定しない。\par

\noindent\textbf{4：代表的条件で適切} 繰返しせん断力。排水が追いつかない飽和土で過剰間隙水圧が上がり、有効応力とせん断抵抗が低下する機構に関係する。\par

\noindent\textbf{5：代表的条件で適切} 乱さない試料採取。密度・粒子骨格等の状態をなるべく保持し、動的試験や液状化抵抗の評価へつなぐ。採取だけで発生の有無を直接決めるものではない。\par

\paragraph{再利用する考え方} 直接測る量と代理指標を区別し、分母・有効応力・水の経路・材料の条件から判断する。

\paragraph{共通理論} \hyperref[theory:engineering:liquefaction]{第\ref*{theory:engineering:liquefaction}節}

\par 参考：\cite{eng52_liquefaction,eng52_plate_liq,geebec_ps}。

\paragraph{関連問題} \hyperref[q:h23-reg-applied-33]{H23-REG-APPLIED-33}

\par 難易度 3（校準初稿）、優先度未評価。

\section{せん断包絡線の切片と傾き}

\subsection{平成29年度 問22}
\label{q:h29-reg-applied-22}

\noindent\textbf{原典：} sources/H29/17\_applied\_science.pdf、PDF物理7頁。

\textbf{公式正答：} 2（single）。

\paragraph{問題の要旨} 一面せん断の直線包絡線から切片と摩擦角を求める。

\Needspace{4\baselineskip}\paragraph{確認した条件・記述（独自要約）}\mbox{}\par

\noindent\textbf{原典照合：} H29応用原典7頁の四測定点・指数4・全五候補を画像確認。\par

\noindent\textbf{条件：} Mohr–Coulombの破壊包絡線を直線とする。指定tan25°=0.47、tan30°=0.58、tan35°=0.70。\par

\noindent\textbf{単位：} σ・τ・cは×10\textsuperscript{4}N/m\textsuperscript{2}。全応力又は有効応力という区別は原典に明示されず、勝手にuの値を追加しない。\par

\paragraph{問われる核心} 一面せん断の直線包絡線から切片と摩擦角を求める。

\paragraph{必要知識} 原典の測定点(σ,τ)は共通単位×10\textsuperscript{4}N/m\textsuperscript{2}で(1,3.1),(2,3.7),(4,4.8),(8,7.1)。端点の傾き(7.1−3.1)/(8−1)=\allowbreak{}0.5714で、指定0.58の30度に近い。φ=30度で各点からc=τ−0.58σを求めると約2.5×10\textsuperscript{4}Paなので2。丸め値は厳密に一つの直線上ではない。最小二乗でも傾き0.5696、切片2.5391となり同じ候補が最も近い。

\paragraph{考え方（査読待ち初稿）} 原典の測定点(σ,τ)は共通単位×10\textsuperscript{4}N/m\textsuperscript{2}で(1, 3.1)、(2, 3.7)、(4, 4.8)、(8, 7.1)。端点の傾き(7.1−3.1)/(8−1)=\allowbreak{}0.5714で、指定0.58の30度に近い。φ=30度で各点からc=τ−0.58σを求めると約2.5×10\textsuperscript{4}Paなので2。丸め値は厳密に一つの直線上ではない。最小二乗でも傾き0.5696、切片2.5391となり同じ候補が最も近い。

\paragraph{正答を選ぶ理由} 原典の測定点(σ,τ)は共通単位×10\textsuperscript{4}N/m\textsuperscript{2}で(1,3.1),(2,3.7),(4,4.8),(8,7.1)。端点の傾き(7.1−3.1)/(8−1)=\allowbreak{}0.5714で、指定0.58の30度に近い。φ=30度で各点からc=τ−0.58σを求めると約2.5×10\textsuperscript{4}Paなので2。丸め値は厳密に一つの直線上ではない。最小二乗でも傾き0.5696、切片2.5391となり同じ候補が最も近い。

\Needspace{4\baselineskip}\paragraph{選択肢の検討（査読待ち初稿）}\mbox{}\par

\noindent\textbf{1：不適切} c=\allowbreak{}2.5×10\textsuperscript{4}Pa、φ=25°。切片は近いが指定tan25°=0.47は測定の傾き約0.57〜0.60から離れる。\par

\noindent\textbf{2：代表的条件で適切} c=\allowbreak{}2.5×10\textsuperscript{4}Pa、φ=30°。指定tan30°=0.58、直線τ=c+σtanφから各点のcは2.52・2.54・2.48・2.46（×10\textsuperscript{4}Pa）で、丸めを含めて2.5と整合。\par

\noindent\textbf{3：不適切} c=\allowbreak{}2.0×10\textsuperscript{4}Pa、φ=25°。切片も傾きも測定と不一致。σ=1で予測τ=2.47に対し3.1。\par

\noindent\textbf{4：不適切} c=\allowbreak{}2.0×10\textsuperscript{4}Pa、φ=30°。傾きは近いがσ=1で予測τ=2.58に対し3.1。切片が不足。\par

\noindent\textbf{5：不適切} c=\allowbreak{}2.0×10\textsuperscript{4}Pa、φ=35°。傾き0.70は大きすぎ、低応力の切片も不足。全測定点で比較する。\par

\paragraph{再利用する考え方} 直接測る量と代理指標を区別し、分母・有効応力・水の経路・材料の条件から判断する。

\paragraph{共通理論} \hyperref[theory:engineering:mohr]{第\ref*{theory:engineering:mohr}節}

\par 参考：\cite{eng52_mohr_ref}。

\par 難易度 3（校準初稿）、優先度未評価。

\section{含水比・間隙率・飽和度の分母}

\subsection{平成25年度 問28}
\label{q:h25-reg-applied-28}

\noindent\textbf{原典：} sources/H25/17\_applied\_science.pdf、PDF物理10頁。

\textbf{公式正答：} 2（single）。

\paragraph{問題の要旨} 含水比・体積含水率・飽和度の分母を区別する。

\Needspace{4\baselineskip}\paragraph{確認した条件・記述（独自要約）}\mbox{}\par

\noindent\textbf{原典照合：} H25原典10頁で全五候補と記号を確認。\par

\noindent\textbf{指定：} 湿潤質量M、全体積V、水質量Mw、水体積Vw、110±5℃炉乾燥質量Ms、間隙体積Vφ。\par

\paragraph{問われる核心} 含水比・体積含水率・飽和度の分母を区別する。

\paragraph{必要知識} 原典指定の炉乾燥質量Msを分母にして水質量Mwを割るので2。水が乾燥土より重ければ含水比は100％を超え得る。分母が湿潤質量なら100％以内となる別量である。乾燥温度110±5℃も原典の定義として保持し、別資料の105℃測定へ置き換えない。

\paragraph{考え方・正答を選ぶ理由（査読待ち初稿）} 原典指定の炉乾燥質量Msを分母にして水質量Mwを割るので2。水が乾燥土より重ければ含水比は100％を超え得る。分母が湿潤質量なら100％以内となる別量である。乾燥温度110±5℃も原典の定義として保持し、別資料の105℃測定へ置き換えない。

\Needspace{4\baselineskip}\paragraph{選択肢の検討（査読待ち初稿）}\mbox{}\par

\noindent\textbf{1：不適切} Mw/M×100。湿潤試料全質量が分母で、乾燥土質量を分母とする土質の含水比とは違う。\par

\noindent\textbf{2：代表的条件で適切} Mw/Ms×100。水の質量を乾燥土粒子の質量で割る含水比。百分率で示すため100を掛ける。\par

\noindent\textbf{3：不適切} Mw/V×100。質量/体積で密度の次元がある。百分率の無次元比ではない。\par

\noindent\textbf{4：不適切} Vw/V×100。全体積に対する水体積の体積含水率。質量の含水比ではない。\par

\noindent\textbf{5：不適切} Vw/Vφ×100。間隙全体積に対する水体積で飽和度。試料全体積の体積含水率とも違う。\par

\paragraph{再利用する考え方} 直接測る量と代理指標を区別し、分母・有効応力・水の経路・材料の条件から判断する。

\paragraph{共通理論} \hyperref[theory:engineering:phases]{第\ref*{theory:engineering:phases}節}

\par 参考：\cite{eng52_soil_phases}。

\paragraph{関連問題} \hyperref[q:h28-reg-applied-24]{H28-REG-APPLIED-24}

\par 難易度 3（校準初稿）、優先度未評価。

\subsection{平成28年度 問24}
\label{q:h28-reg-applied-24}

\noindent\textbf{原典：} sources/H28/17\_applied\_science.pdf、PDF物理10頁。

\textbf{公式正答：} 3（single）。

\paragraph{問題の要旨} 三相の体積から間隙率と体積含水率を組み立てる。

\Needspace{4\baselineskip}\paragraph{確認した条件・記述（独自要約）}\mbox{}\par

\noindent\textbf{原典照合：} H28原典10頁で組合せ表全行を確認。\par

\noindent\textbf{記号：} Vs固相、Vw水、Va空気、V=\allowbreak{}Vs+Vw+Va。候補の百分率係数100は共通。\par

\paragraph{問われる核心} 三相の体積から間隙率と体積含水率を組み立てる。

\paragraph{必要知識} 全体積V=\allowbreak{}Vs+Vw+Vaと置く。間隙率n=\allowbreak{}(Vw+Va)/V、体積含水率θ=Vw/Vで、それぞれ百分率なら100を掛ける。両方の分母が全体積となる3を選ぶ。間隙比eと飽和度Srは別分母である。

\paragraph{考え方・正答を選ぶ理由（査読待ち初稿）} 全体積V=\allowbreak{}Vs+Vw+Vaと置く。間隙率n=\allowbreak{}(Vw+Va)/V、体積含水率θ=Vw/Vで、それぞれ百分率なら100を掛ける。両方の分母が全体積となる3を選ぶ。間隙比eと飽和度Srは別分母である。

\Needspace{4\baselineskip}\paragraph{選択肢の検討（査読待ち初稿）}\mbox{}\par

\noindent\textbf{1：不適切} A=\allowbreak{}(Vw+Va)/V、B=\allowbreak{}Vw/Vs。Aは間隙率だがBの分母が固相で、全体積の体積含水率ではない。\par

\noindent\textbf{2：不適切} A=\allowbreak{}(Vw+Va)/V、B=\allowbreak{}Vw/(Vw+Va)。Bは飽和度。\par

\noindent\textbf{3：代表的条件で適切} A=\allowbreak{}(Vw+Va)/V、B=\allowbreak{}Vw/V。間隙率も体積含水率も全体積V=\allowbreak{}Vs+Vw+Vaが分母。Aの分子は水と空気、Bの分子は水。\par

\noindent\textbf{4：不適切} A=\allowbreak{}(Vw+Va)/Vs、B=\allowbreak{}Vw/(Vw+Va)。Aは間隙比、Bは飽和度で両方違う。\par

\noindent\textbf{5：不適切} A=\allowbreak{}(Vw+Va)/Vs、B=\allowbreak{}Vw/V。Bは正しいがAは間隙比。\par

\paragraph{再利用する考え方} 直接測る量と代理指標を区別し、分母・有効応力・水の経路・材料の条件から判断する。

\paragraph{共通理論} \hyperref[theory:engineering:phases]{第\ref*{theory:engineering:phases}節}

\par 参考：\cite{eng52_soil_phases}。

\paragraph{関連問題} \hyperref[q:h25-reg-applied-28]{H25-REG-APPLIED-28}

\par 難易度 3（校準初稿）、優先度未評価。

\section{黄鉄鉱の酸化・溶解・塩析出}

\subsection{平成24年度 問30}
\label{q:h24-reg-applied-30}

\noindent\textbf{原典：} sources/H24/17\_applied\_science.pdf、PDF物理9頁。

\textbf{公式正答：} 5（single）。

\par\textbf{検討未確定：} 酸化フロントにすべりが多いという頻度、泥岩の全岩種・全層序への一般化は未確定。未酸化試料の掘削後酸化と硫酸塩析出を時間順で補足。

\paragraph{問題の要旨} 海成泥岩の黄鉄鉱酸化・酸性排水・析出を区別する。

\Needspace{4\baselineskip}\paragraph{確認した条件・記述（独自要約）}\mbox{}\par

\noindent\textbf{原典照合：} 22原典画像で25問の全五肢・数値・単位・否定を確認。\par

\noindent\textbf{対象：} 海成泥岩の黄鉄鉱酸化・酸性排水・析出を区別する。\par

\paragraph{問われる核心} 海成泥岩の黄鉄鉱酸化・酸性排水・析出を区別する。

\paragraph{必要知識} 誤り5を、黄鉄鉱の酸化が生む酸と弱化の過程から説明する。生成した硫酸イオンがCaと結び石膏等で析出すると盤膨れの別過程になり得る。層理・凝灰岩の弱層・地下水経路を合わせ、岩種と酸化だけで全地点のすべりを決めない。

\paragraph{考え方・正答を選ぶ理由（査読待ち初稿）} 誤り5を、黄鉄鉱の酸化が生む酸と弱化の過程から説明する。生成した硫酸イオンがCaと結び石膏等で析出すると盤膨れの別過程になり得る。層理・凝灰岩の弱層・地下水経路を合わせ、岩種と酸化だけで全地点のすべりを決めない。

\Needspace{4\baselineskip}\paragraph{選択肢の検討（査読待ち初稿）}\mbox{}\par

\noindent\textbf{1：代表的条件で適切} すべり面が層理に平行。弱い層・連続する面の向きがすべりの経路に関係する。全すべり面が層理という限定ではない。\par

\noindent\textbf{2：代表的条件で適切} 凝灰岩にすべり面。砂岩泥岩の間にある変質・軟化しやすい凝灰岩等が弱層となることがある。\par

\noindent\textbf{3：条件・全称を留保} 酸化フロント・下の溶解帯。黄鉄鉱の酸化と酸性水による膠結物の溶解は弱化へつながる。原典の多いという頻度と特定のすべり位置の全称は追加照合対象。\par

\noindent\textbf{4：代表的条件で適切} 未酸化青黒泥岩と硫酸塩・石膏の盤膨れ。新鮮な還元状態の岩石が掘削後に水・酸素と反応し、硫酸塩が析出して結晶成長圧が生じ得る。未酸化のままで酸化反応が進むという矛盾した理解は避ける。\par

\noindent\textbf{5：不適切} 黄鉄鉱でアルカリ性排水。黄鉄鉱の酸化は酸を生む。炭酸塩等で中和されることはあっても、その酸化自体をアルカリ化の起源とはしない。\par

\paragraph{再利用する考え方} 直接測る量と代理指標を区別し、分母・有効応力・水の経路・材料の条件から判断する。

\paragraph{共通理論} \hyperref[theory:engineering:pyrite]{第\ref*{theory:engineering:pyrite}節}

\par 参考：\cite{eng52_weathering,eng52_slip_weaklayer}。

\paragraph{関連問題} \hyperref[q:r01-reg-applied-29]{R01-REG-APPLIED-29}

\par 難易度 3（校準初稿）、優先度未評価。

\subsection{令和元年度 問29}
\label{q:r01-reg-applied-29}

\noindent\textbf{原典：} sources/R01/17\_applied\_science.pdf、PDF物理10頁。

\textbf{公式正答：} 5（single）。

\par\textbf{検討未確定：} 酸化フロントにすべりが多いという頻度、泥岩の全岩種・全層序への一般化は未確定。未酸化試料の掘削後酸化と硫酸塩析出を時間順で補足。

\paragraph{問題の要旨} 海成泥岩の黄鉄鉱酸化・酸性排水・析出を区別する。

\Needspace{4\baselineskip}\paragraph{確認した条件・記述（独自要約）}\mbox{}\par

\noindent\textbf{原典照合：} 22原典画像で25問の全五肢・数値・単位・否定を確認。\par

\noindent\textbf{対象：} 海成泥岩の黄鉄鉱酸化・酸性排水・析出を区別する。\par

\paragraph{問われる核心} 海成泥岩の黄鉄鉱酸化・酸性排水・析出を区別する。

\paragraph{必要知識} 誤り5を、黄鉄鉱の酸化が生む酸と弱化の過程から説明する。生成した硫酸イオンがCaと結び石膏等で析出すると盤膨れの別過程になり得る。層理・凝灰岩の弱層・地下水経路を合わせ、岩種と酸化だけで全地点のすべりを決めない。

\paragraph{考え方・正答を選ぶ理由（査読待ち初稿）} 誤り5を、黄鉄鉱の酸化が生む酸と弱化の過程から説明する。生成した硫酸イオンがCaと結び石膏等で析出すると盤膨れの別過程になり得る。層理・凝灰岩の弱層・地下水経路を合わせ、岩種と酸化だけで全地点のすべりを決めない。

\Needspace{4\baselineskip}\paragraph{選択肢の検討（査読待ち初稿）}\mbox{}\par

\noindent\textbf{1：代表的条件で適切} 未酸化青黒泥岩と硫酸塩・石膏の盤膨れ。新鮮な還元状態の岩石が掘削後に水・酸素と反応し、硫酸塩が析出して結晶成長圧が生じ得る。未酸化のままで酸化反応が進むという矛盾した理解は避ける。\par

\noindent\textbf{2：代表的条件で適切} 黄鉄鉱で酸性排水の可能性。酸素・水と反応すると硫酸イオンと酸を生む。反応量・緩衝能で実際の排水pHは変わる。\par

\noindent\textbf{3：代表的条件で適切} 凝灰岩にすべり面。砂岩泥岩の間にある変質・軟化しやすい凝灰岩等が弱層となることがある。\par

\noindent\textbf{4：代表的条件で適切} すべり面が層理に平行。弱い層・連続する面の向きがすべりの経路に関係する。全すべり面が層理という限定ではない。\par

\noindent\textbf{5：不適切} 酸化フロントが地すべりを抑制。酸化・溶解風化で結合が弱まり、すべりへつながり得るため、一方向の抑制とはしない。酸化生成物の間隙充填等がある場合も条件を分ける。\par

\paragraph{再利用する考え方} 直接測る量と代理指標を区別し、分母・有効応力・水の経路・材料の条件から判断する。

\paragraph{共通理論} \hyperref[theory:engineering:pyrite]{第\ref*{theory:engineering:pyrite}節}

\par 参考：\cite{eng52_weathering,eng52_slip_weaklayer}。

\paragraph{関連問題} \hyperref[q:h24-reg-applied-30]{H24-REG-APPLIED-30}

\par 難易度 3（校準初稿）、優先度未評価。

\section{地すべり・崩壊の代表比較}

\subsection{平成30年度 問24}
\label{q:h30-reg-applied-24}

\noindent\textbf{原典：} sources/H30/17\_applied\_science.pdf、PDF物理8頁。

\textbf{公式正答：} 3（single）。

\par\textbf{検討未確定：} 原典の塑性/脆性、緩慢/高速、砂質/粘性の全材料への排他性は未確認。公式選択と実際の多様な移動機構を分けて要査読とする。

\paragraph{問題の要旨} 地すべりと崩壊の代表比較を全称にしない。

\Needspace{4\baselineskip}\paragraph{確認した条件・記述（独自要約）}\mbox{}\par

\noindent\textbf{原典照合：} 22原典画像で25問の全五肢・数値・単位・否定を確認。\par

\noindent\textbf{対象：} 地すべりと崩壊の代表比較を全称にしない。\par

\paragraph{問われる核心} 地すべりと崩壊の代表比較を全称にしない。

\paragraph{必要知識} 公式3は砂質・粘性土の典型対応を逆にした候補。移動の継続性、弱面、水圧、降雨・地震、兆候の観察を合わせる。速度と破壊形式は代表比較であり、全事例を二種類へ固定する定義ではない。

\paragraph{考え方・正答を選ぶ理由（査読待ち初稿）} 公式3は砂質・粘性土の典型対応を逆にした候補。移動の継続性、弱面、水圧、降雨・地震、兆候の観察を合わせる。速度と破壊形式は代表比較であり、全事例を二種類へ固定する定義ではない。

\Needspace{4\baselineskip}\paragraph{選択肢の検討（査読待ち初稿）}\mbox{}\par

\noindent\textbf{1：条件・全称を留保} 地すべり塑性・崩壊脆性。原典の典型比較として保持するが、全地すべりを塑性・全崩壊を脆性という二分法で定義しない。せん断面の滑動と岩塊の移動には多様な機構がある。\par

\noindent\textbf{2：条件・全称を留保} 地すべり緩慢・崩壊高速。代表的な観察上の比較。地すべりにも急激な滑動があるため、名称から全速度を固定しない。\par

\noindent\textbf{3：不適切} 地すべり砂質・崩壊粘性。通常の教材の代表比較では、地すべりは粘性土等の弱いすべり面、崩壊は表層の砂質土等と関係し、この対応は逆。全材料を排他的に分類する規則ではない。\par

\noindent\textbf{4：代表的条件で適切} 地すべり地下水・崩壊豪雨地震。地下水圧や雨・地震が安定を低下させる。地すべりにも豪雨・地震が働き、唯一の引き金を分けるものではない。\par

\noindent\textbf{5：代表的条件で適切} 地すべりに亀裂・隆起・陥没の兆候。長期滑動等では変形の兆候を認める場合がある。崩壊に兆候が全くない、地すべりなら必ず発見できるとはしない。\par

\paragraph{再利用する考え方} 直接測る量と代理指標を区別し、分母・有効応力・水の経路・材料の条件から判断する。

\paragraph{共通理論} \hyperref[theory:engineering:landslide]{第\ref*{theory:engineering:landslide}節}

\par 参考：\cite{eng52_landslide_ref}。

\paragraph{関連問題} \hyperref[q:r04-reg-applied-27]{R04-REG-APPLIED-27}

\par 難易度 3（校準初稿）、優先度未評価。

\subsection{令和4年度 問27}
\label{q:r04-reg-applied-27}

\noindent\textbf{原典：} sources/R04/17\_applied\_science.pdf、PDF物理9頁。

\textbf{公式正答：} 4（single）。

\par\textbf{検討未確定：} 原典の塑性/脆性、緩慢/高速、砂質/粘性の全材料への排他性は未確認。公式選択と実際の多様な移動機構を分けて要査読とする。

\paragraph{問題の要旨} 地すべりと崩壊の代表比較を全称にしない。

\Needspace{4\baselineskip}\paragraph{確認した条件・記述（独自要約）}\mbox{}\par

\noindent\textbf{原典照合：} 22原典画像で25問の全五肢・数値・単位・否定を確認。\par

\noindent\textbf{対象：} 地すべりと崩壊の代表比較を全称にしない。\par

\paragraph{問われる核心} 地すべりと崩壊の代表比較を全称にしない。

\paragraph{必要知識} 公式4は砂質・粘性土の典型対応を逆にした候補。移動の継続性、弱面、水圧、降雨・地震、兆候の観察を合わせる。速度と破壊形式は代表比較であり、全事例を二種類へ固定する定義ではない。

\paragraph{考え方・正答を選ぶ理由（査読待ち初稿）} 公式4は砂質・粘性土の典型対応を逆にした候補。移動の継続性、弱面、水圧、降雨・地震、兆候の観察を合わせる。速度と破壊形式は代表比較であり、全事例を二種類へ固定する定義ではない。

\Needspace{4\baselineskip}\paragraph{選択肢の検討（査読待ち初稿）}\mbox{}\par

\noindent\textbf{1：条件・全称を留保} 地すべり塑性・崩壊脆性。原典の典型比較として保持するが、全地すべりを塑性・全崩壊を脆性という二分法で定義しない。せん断面の滑動と岩塊の移動には多様な機構がある。\par

\noindent\textbf{2：代表的条件で適切} 地すべりに亀裂・隆起・陥没の兆候。長期滑動等では変形の兆候を認める場合がある。崩壊に兆候が全くない、地すべりなら必ず発見できるとはしない。\par

\noindent\textbf{3：代表的条件で適切} 地すべり地下水・崩壊豪雨地震。地下水圧や雨・地震が安定を低下させる。地すべりにも豪雨・地震が働き、唯一の引き金を分けるものではない。\par

\noindent\textbf{4：不適切} 地すべり砂質・崩壊粘性。通常の教材の代表比較では、地すべりは粘性土等の弱いすべり面、崩壊は表層の砂質土等と関係し、この対応は逆。全材料を排他的に分類する規則ではない。\par

\noindent\textbf{5：条件・全称を留保} 地すべり緩慢・崩壊高速。代表的な観察上の比較。地すべりにも急激な滑動があるため、名称から全速度を固定しない。\par

\paragraph{再利用する考え方} 直接測る量と代理指標を区別し、分母・有効応力・水の経路・材料の条件から判断する。

\paragraph{共通理論} \hyperref[theory:engineering:landslide]{第\ref*{theory:engineering:landslide}節}

\par 参考：\cite{eng52_landslide_ref}。

\paragraph{関連問題} \hyperref[q:h30-reg-applied-24]{H30-REG-APPLIED-24}

\par 難易度 3（校準初稿）、優先度未評価。

\section{せん断強度・貫入と圧密の測定量}

\subsection{平成23年度 問34}
\label{q:h23-reg-applied-34}

\noindent\textbf{原典：} sources/H23/17\_applied\_science.pdf、PDF物理11頁。

\textbf{公式正答：} 3（single）。

\paragraph{問題の要旨} 圧密の体積変形とせん断強度試験の測定量を分ける。

\Needspace{4\baselineskip}\paragraph{確認した条件・記述（独自要約）}\mbox{}\par

\noindent\textbf{原典照合：} 22原典画像で25問の全五肢・数値・単位・否定を確認。\par

\noindent\textbf{対象：} 圧密の体積変形とせん断強度試験の測定量を分ける。\par

\paragraph{問われる核心} 圧密の体積変形とせん断強度試験の測定量を分ける。

\paragraph{必要知識} 最も関係の薄い候補3は圧密試験。せん断破壊、回転トルク、貫入抵抗を強度・支持性へ結び付ける他の試験と、圧密の圧縮・排水・時間の主目的を分ける。圧密と強度の状態依存まで無関係とはしない。

\paragraph{考え方・正答を選ぶ理由（査読待ち初稿）} 最も関係の薄い候補3は圧密試験。せん断破壊、回転トルク、貫入抵抗を強度・支持性へ結び付ける他の試験と、圧密の圧縮・排水・時間の主目的を分ける。圧密と強度の状態依存まで無関係とはしない。

\Needspace{4\baselineskip}\paragraph{選択肢の検討（査読待ち初稿）}\mbox{}\par

\noindent\textbf{1：代表的条件で適切} 三軸又は一軸の力学試験。拘束圧・排水条件を指定して試料を破壊へ載荷し、強度や応力ひずみ応答を調べる。\par

\noindent\textbf{2：代表的条件で適切} ベーンせん断試験。羽根を回すトルク等から軟らかい粘性土の非排水せん断抵抗を評価する。\par

\noindent\textbf{3：不適切} 圧密試験。主目的は荷重増加に対する時間依存の体積変化・圧縮性・排水の評価。せん断破壊強度を直接測る試験とは違う。\par

\noindent\textbf{4：代表的条件で適切} スウェーデン式サウンディング。貫入抵抗から地盤の硬軟や支持性・土質を経験的に推定する。推定量と直接破壊試験を区別する。\par

\noindent\textbf{5：代表的条件で適切} 標準貫入試験。N値は所定貫入の打撃数で、密度・硬軟等の経験的評価に使う。N値をせん断強度の単位としない。\par

\paragraph{再利用する考え方} 直接測る量と代理指標を区別し、分母・有効応力・水の経路・材料の条件から判断する。

\paragraph{共通理論} \hyperref[theory:engineering:tests]{第\ref*{theory:engineering:tests}節}

\par 参考：\cite{eng52_geotech_tests,eng52_consolidation}。

\paragraph{関連問題} \hyperref[q:h26-reg-applied-34]{H26-REG-APPLIED-34}

\par 難易度 2（校準初稿）、優先度未評価。

\subsection{平成26年度 問34}
\label{q:h26-reg-applied-34}

\noindent\textbf{原典：} sources/H26/17\_applied\_science.pdf、PDF物理14頁。

\textbf{公式正答：} 1（single）。

\paragraph{問題の要旨} 圧密の体積変形とせん断強度試験の測定量を分ける。

\Needspace{4\baselineskip}\paragraph{確認した条件・記述（独自要約）}\mbox{}\par

\noindent\textbf{原典照合：} 22原典画像で25問の全五肢・数値・単位・否定を確認。\par

\noindent\textbf{対象：} 圧密の体積変形とせん断強度試験の測定量を分ける。\par

\paragraph{問われる核心} 圧密の体積変形とせん断強度試験の測定量を分ける。

\paragraph{必要知識} 最も関係の薄い候補1は圧密試験。せん断破壊、回転トルク、貫入抵抗を強度・支持性へ結び付ける他の試験と、圧密の圧縮・排水・時間の主目的を分ける。圧密と強度の状態依存まで無関係とはしない。

\paragraph{考え方・正答を選ぶ理由（査読待ち初稿）} 最も関係の薄い候補1は圧密試験。せん断破壊、回転トルク、貫入抵抗を強度・支持性へ結び付ける他の試験と、圧密の圧縮・排水・時間の主目的を分ける。圧密と強度の状態依存まで無関係とはしない。

\Needspace{4\baselineskip}\paragraph{選択肢の検討（査読待ち初稿）}\mbox{}\par

\noindent\textbf{1：不適切} 圧密試験。主目的は荷重増加に対する時間依存の体積変化・圧縮性・排水の評価。せん断破壊強度を直接測る試験とは違う。\par

\noindent\textbf{2：代表的条件で適切} 三軸又は一軸の力学試験。拘束圧・排水条件を指定して試料を破壊へ載荷し、強度や応力ひずみ応答を調べる。\par

\noindent\textbf{3：代表的条件で適切} 標準貫入試験。N値は所定貫入の打撃数で、密度・硬軟等の経験的評価に使う。N値をせん断強度の単位としない。\par

\noindent\textbf{4：代表的条件で適切} ベーンせん断試験。羽根を回すトルク等から軟らかい粘性土の非排水せん断抵抗を評価する。\par

\noindent\textbf{5：代表的条件で適切} スウェーデン式サウンディング。貫入抵抗から地盤の硬軟や支持性・土質を経験的に推定する。推定量と直接破壊試験を区別する。\par

\paragraph{再利用する考え方} 直接測る量と代理指標を区別し、分母・有効応力・水の経路・材料の条件から判断する。

\paragraph{共通理論} \hyperref[theory:engineering:tests]{第\ref*{theory:engineering:tests}節}

\par 参考：\cite{eng52_geotech_tests,eng52_consolidation}。

\paragraph{関連問題} \hyperref[q:h23-reg-applied-34]{H23-REG-APPLIED-34}, \hyperref[q:h30-reg-applied-27]{H30-REG-APPLIED-27}

\par 難易度 2（校準初稿）、優先度未評価。

\subsection{平成30年度 問27}
\label{q:h30-reg-applied-27}

\noindent\textbf{原典：} sources/H30/17\_applied\_science.pdf、PDF物理9頁。

\textbf{公式正答：} 1（single）。

\paragraph{問題の要旨} 圧密の体積変形とせん断強度試験の測定量を分ける。

\Needspace{4\baselineskip}\paragraph{確認した条件・記述（独自要約）}\mbox{}\par

\noindent\textbf{原典照合：} 22原典画像で25問の全五肢・数値・単位・否定を確認。\par

\noindent\textbf{対象：} 圧密の体積変形とせん断強度試験の測定量を分ける。\par

\paragraph{問われる核心} 圧密の体積変形とせん断強度試験の測定量を分ける。

\paragraph{必要知識} 最も関係の薄い候補1は圧密試験。せん断破壊、回転トルク、貫入抵抗を強度・支持性へ結び付ける他の試験と、圧密の圧縮・排水・時間の主目的を分ける。圧密と強度の状態依存まで無関係とはしない。

\paragraph{考え方・正答を選ぶ理由（査読待ち初稿）} 最も関係の薄い候補1は圧密試験。せん断破壊、回転トルク、貫入抵抗を強度・支持性へ結び付ける他の試験と、圧密の圧縮・排水・時間の主目的を分ける。圧密と強度の状態依存まで無関係とはしない。

\Needspace{4\baselineskip}\paragraph{選択肢の検討（査読待ち初稿）}\mbox{}\par

\noindent\textbf{1：不適切} 圧密試験。主目的は荷重増加に対する時間依存の体積変化・圧縮性・排水の評価。せん断破壊強度を直接測る試験とは違う。\par

\noindent\textbf{2：代表的条件で適切} 三軸又は一軸の力学試験。拘束圧・排水条件を指定して試料を破壊へ載荷し、強度や応力ひずみ応答を調べる。\par

\noindent\textbf{3：代表的条件で適切} 標準貫入試験。N値は所定貫入の打撃数で、密度・硬軟等の経験的評価に使う。N値をせん断強度の単位としない。\par

\noindent\textbf{4：代表的条件で適切} ベーンせん断試験。羽根を回すトルク等から軟らかい粘性土の非排水せん断抵抗を評価する。\par

\noindent\textbf{5：代表的条件で適切} スウェーデン式サウンディング。貫入抵抗から地盤の硬軟や支持性・土質を経験的に推定する。推定量と直接破壊試験を区別する。\par

\paragraph{再利用する考え方} 直接測る量と代理指標を区別し、分母・有効応力・水の経路・材料の条件から判断する。

\paragraph{共通理論} \hyperref[theory:engineering:tests]{第\ref*{theory:engineering:tests}節}

\par 参考：\cite{eng52_geotech_tests,eng52_consolidation}。

\paragraph{関連問題} \hyperref[q:h26-reg-applied-34]{H26-REG-APPLIED-34}

\par 難易度 2（校準初稿）、優先度未評価。

\section{揚水水頭低下と粘土圧密沈下}

\subsection{平成25年度 問30}
\label{q:h25-reg-applied-30}

\noindent\textbf{原典：} sources/H25/17\_applied\_science.pdf、PDF物理11頁。

\textbf{公式正答：} 3（single）。

\paragraph{問題の要旨} 揚水による水圧低下と粘土層の圧密沈下をつなぐ。

\Needspace{4\baselineskip}\paragraph{確認した条件・記述（独自要約）}\mbox{}\par

\noindent\textbf{原典照合：} 22原典画像で25問の全五肢・数値・単位・否定を確認。\par

\noindent\textbf{対象：} 揚水による水圧低下と粘土層の圧密沈下をつなぐ。\par

\paragraph{問われる核心} 揚水による水圧低下と粘土層の圧密沈下をつなぐ。

\paragraph{必要知識} 3は不圧水位へ主因を限定した点が不適切。σ\ensuremath{'}=σ−uで、水圧uが下がると骨格の有効応力が増し、圧縮性の高い粘土層の排水・収縮が地表へ現れる。層ごとの水頭と地盤高を観測して時間差も追う。

\paragraph{考え方・正答を選ぶ理由（査読待ち初稿）} 3は不圧水位へ主因を限定した点が不適切。σ\ensuremath{'}=σ−uで、水圧uが下がると骨格の有効応力が増し、圧縮性の高い粘土層の排水・収縮が地表へ現れる。層ごとの水頭と地盤高を観測して時間差も追う。

\Needspace{4\baselineskip}\paragraph{選択肢の検討（査読待ち初稿）}\mbox{}\par

\noindent\textbf{1：代表的条件で適切} 構造物浮き上がり・道路変形・埋設破損。周囲の地盤が沈むと杭等に支えられた構造物が相対的に浮き上がった状態となり、道路・埋設物へ被害が生じる。\par

\noindent\textbf{2：代表的条件で適切} 帯水層に接する粘土も圧密収縮。揚水で水頭・間隙水圧が低下し、同じ全応力のもとで土粒子骨格の有効応力が増加する。\par

\noindent\textbf{3：不適切} 主に不圧地下水位の変動。広域の揚水沈下では被圧帯水層の水頭低下と隣接する粘土層の圧密が重要。不圧水位だけへ限定しない。\par

\noindent\textbf{4：代表的条件で適切} 低平地の浸水・冠水危険増大。標高が低下し排水勾配や治水の条件が悪化し、二次的被害のリスクを高める。\par

\noindent\textbf{5：代表的条件で適切} 帯水層深度ごとの観測井。水頭が異なる層を区別して観測する。複数帯水層をつなぐ井戸の混合水位だけでは各層を評価できない。\par

\paragraph{再利用する考え方} 直接測る量と代理指標を区別し、分母・有効応力・水の経路・材料の条件から判断する。

\paragraph{共通理論} \hyperref[theory:engineering:consolidation]{第\ref*{theory:engineering:consolidation}節}

\par 参考：\cite{eng52_subsidence_ref,eng52_consolidation}。

\par 難易度 2（校準初稿）、優先度未評価。

\section{切羽・偏圧・山はねと地表沈下}

\subsection{令和3年度 問28}
\label{q:r03-reg-applied-28}

\noindent\textbf{原典：} sources/R03/17\_applied\_science.pdf、PDF物理10頁。

\textbf{公式正答：} 2（single）。

\par\textbf{検討未確定：} 緩み土圧と応力損傷の全条件の境界、土かぶりと地表沈下の定量関係は未確定。代表傾向として保持し全工法へ一般化しない。

\paragraph{問題の要旨} 切羽崩壊・偏圧・山はねと土かぶりを区別する。

\Needspace{4\baselineskip}\paragraph{確認した条件・記述（独自要約）}\mbox{}\par

\noindent\textbf{原典照合：} 22原典画像で25問の全五肢・数値・単位・否定を確認。\par

\noindent\textbf{対象：} 切羽崩壊・偏圧・山はねと土かぶりを区別する。\par

\paragraph{問われる核心} 切羽崩壊・偏圧・山はねと土かぶりを区別する。

\paragraph{必要知識} 適切2は切羽の構造・水と不安定化の対応。応力の解放・集中、高応力の硬岩の山はね、割れ目に沿う緩み、非対称な偏圧、地表沈下の各機構を別々に比較する。

\paragraph{考え方・正答を選ぶ理由（査読待ち初稿）} 適切2は切羽の構造・水と不安定化の対応。応力の解放・集中、高応力の硬岩の山はね、割れ目に沿う緩み、非対称な偏圧、地表沈下の各機構を別々に比較する。

\Needspace{4\baselineskip}\paragraph{選択肢の検討（査読待ち初稿）}\mbox{}\par

\noindent\textbf{1：不適切} 割れ目の少ない硬岩の崩壊を緩み土圧。高応力の健全硬岩での損傷・山はねと、割れ目のある地山の重力的な緩みを混同している。現象名を応力・不連続面から分ける。\par

\noindent\textbf{2：代表的条件で適切} 切羽天端・鏡面の構造と湧水による崩壊。不連続面や弱い層、水による切羽の不安定化を観察する。施工中の観測を支保・設計へ反映する。\par

\noindent\textbf{3：不適切} 緩傾斜・水平地層で偏圧。典型的な偏圧は斜面地形や傾斜地層等の非対称な条件で問題になる。水平・対称をその原因へ固定しない。\par

\noindent\textbf{4：不適切} 小土かぶり・多節理で山はね。代表的な山はねは高応力・大土かぶり・硬く節理の少ない岩盤で起こりやすい。原典は条件を逆にした。\par

\noindent\textbf{5：条件・全称を留保} 土かぶりが大きいほど地表沈下大。地表沈下は浅い地山・緩み・水位低下・地層の条件等に依存する。大土かぶりだけで地表沈下が大きくなるという単純な単調関係にしない。\par

\paragraph{再利用する考え方} 直接測る量と代理指標を区別し、分母・有効応力・水の経路・材料の条件から判断する。

\paragraph{共通理論} \hyperref[theory:engineering:tunnel]{第\ref*{theory:engineering:tunnel}節}

\par 参考：\cite{eng52_tunnel_ref,eng52_rockburst}。

\par 難易度 3（校準初稿）、優先度未評価。

\section{岩石試験・岩盤強度の代理指標}

\subsection{令和元年度再試験 問27}
\label{q:r01-ret-applied-27}

\noindent\textbf{原典：} sources/R01\_retest/17\_applied\_science.pdf、PDF物理9頁。

\textbf{公式正答：} 3（single）。

\par\textbf{検討未確定：} 岩盤評価の相対比較である。RQD・打診・波速から強度を一意に決めず、弱岩の圧縮降伏の例外を保持。

\paragraph{問題の要旨} 岩盤強度の直接試験と間接指標、圧密降伏を比較する。

\Needspace{4\baselineskip}\paragraph{確認した条件・記述（独自要約）}\mbox{}\par

\noindent\textbf{原典照合：} 29問21原典画像で否定、全五候補、数式・上付き単位・観測対象を照合。\par

\noindent\textbf{前提：} 岩石試験片の破壊強度と、割れ目を含む岩盤の状態及び土の圧密履歴を区別する。\par

\paragraph{問われる核心} 岩盤強度の直接試験と間接指標、圧密降伏を比較する。

\paragraph{必要知識} 岩石試験片の破壊強度と、割れ目を含む岩盤の状態及び土の圧密履歴を区別する。

\paragraph{考え方・正答を選ぶ理由（査読待ち初稿）} 3の圧密降伏応力である。通常は土の一次元圧縮履歴と降伏を表す量であり、他の項目は岩片強度・割れ目状態・弾性応答など岩盤評価へ直接又は間接に使う。ただし弱い泥岩等の圧縮降伏もあり、関連性を完全ゼロとはしない。

\Needspace{4\baselineskip}\paragraph{選択肢の検討（査読待ち初稿）}\mbox{}\par

\noindent\textbf{1：条件付きで適切} ハンマー打診音。硬さ・風化・割れ目等の現地判別に用いるが、定量強度の直接測定ではなく観察者と条件に依存する。\par

\noindent\textbf{2：適切} 一軸圧縮強度。試験片を側方無拘束で圧縮破壊させた最大応力であり、岩片材料の強度を直接評価する。岩盤の弱面は別である。\par

\noindent\textbf{3：不適切} 圧密降伏応力。通常の土の圧密・応力履歴の指標であり、岩盤強度特性の比較では相対的に関係が薄い。弱岩の圧縮降伏との関連は否定しない。\par

\noindent\textbf{4：条件付きで適切} RQD。健全コアの所定長さ割合により割れ目による岩盤状態を表す。長さ品質指標そのものが一軸強度ではない。\par

\noindent\textbf{5：条件付きで適切} 弾性波速度。弾性率、密度、割れ目、含水と関係し岩盤評価に使う。破壊強度との一意換算ではない。\par

\paragraph{再利用する考え方} 岩片の材料試験と岩盤の割れ目・水・応力条件、代理指標を分ける。

\paragraph{共通理論} \hyperref[theory:engineering:strengthindices]{第\ref*{theory:engineering:strengthindices}節}

\par 参考：\cite{jseg_engineering_faq,geebec_ps}。

\par 難易度 3（校準初稿）、優先度未評価。

\section{岩石風化・変質・スレーキング}

\subsection{平成29年度 問24}
\label{q:h29-reg-applied-24}

\noindent\textbf{原典：} sources/H29/17\_applied\_science.pdf、PDF物理8頁。

\textbf{公式正答：} 3（single）。

\paragraph{問題の要旨} 風化・熱水変質・スレーキングが岩石強度へ及ぼす作用を判定する。

\Needspace{4\baselineskip}\paragraph{確認した条件・記述（独自要約）}\mbox{}\par

\noindent\textbf{方向：} 最も不適切なもの。\par

\noindent\textbf{対象：} 一般的な岩石分類・形成過程。\par

\paragraph{問われる核心} 風化・変質と間隙内圧による骨格破壊

\paragraph{必要知識} 風化・変質と間隙内圧による骨格破壊。母材・運搬堆積・有機物起源、風化、弱面、水の作用、岩盤と土の物性。

\paragraph{考え方・正答を選ぶ理由（査読待ち初稿）} スレーキングは水への浸漬などで岩塊が崩れる現象である。封入空気モデルでは、水の侵入で圧縮された空気の内圧が骨格を引っ張るので、岩塊の圧縮力とした3が不適切である。空気の状態変化と固体の応力を別々に記述する。膨潤・吸着・化学劣化なども機構となり得る。

\Needspace{4\baselineskip}\paragraph{選択肢の検討（査読待ち初稿）}\mbox{}\par

\noindent\textbf{1：適切（一般的分類）} 有色鉱物に富む黒い火成岩ほど一般に風化に不安定である。。高温形成のFe・Mgに富む鉱物は地表環境で変化しやすい。色そのものを原因とする法則ではなく、割れ目や組織も耐風化性を左右する。\par

\noindent\textbf{2：適切（一般的分類）} 磁鉄鉱から赤鉄鉱・褐鉄鉱への酸化が風化層の色に現れる。。Feの酸化と水和等により赤褐色系の鉱物を生じる。褐鉄鉱は単一の理想化学式を持つ一鉱物としてではなく含水酸化鉄の集合の名称として扱う。\par

\noindent\textbf{3：不適切} 浸水で間隙空気が圧縮され、岩塊内に圧縮力を生じてスレーキングする。。空気は圧縮されるが、内圧が孔壁・骨格を押し広げて引張破壊へ働く。空気の圧縮と岩塊への圧縮を混同している。封入空気だけが全試料の機構とも限らない。\par

\noindent\textbf{4：適切（一般的分類）} グリーンタフ地域の新第三紀層は熱水変質で軟弱化していることが多い。。火山性の岩石・堆積物が熱水反応で粘土化等を受け、低強度となる場合がある。全てのグリーンタフが同じ弱さという意味ではない。\par

\noindent\textbf{5：適切（一般的分類）} 一般に風化・変質は弱化させるが、珪化は強度を増すこともある。。シリカの充填・交代は緻密化や結合強化をもたらす場合があり、変質=常に強度低下とはしない。\par

\paragraph{再利用する考え方} 風化・変質と間隙内圧による骨格破壊を成因と観察される性質に結び付け、例示と全称命題、主要鉱物と副成分、公式解答と未解決の解釈を区別する。

\paragraph{共通理論} \hyperref[theory:engineering:weathering]{第\ref*{theory:engineering:weathering}節}, \hyperref[theory:engineering:slaking]{第\ref*{theory:engineering:slaking}節}

\par 参考：\cite{slaking1979,slaking2006,jseg_engineering_faq}。

\par 難易度 3（校準初稿）、優先度未評価。

\section{特殊土の母材・成因}

\subsection{平成25年度 問19}
\label{q:h25-reg-applied-19}

\noindent\textbf{原典：} sources/H25/17\_applied\_science.pdf、PDF物理8頁。

\textbf{公式正答：} 5（single）。

\paragraph{問題の要旨} まさ土・しらす・黄土・火山灰質粘性土・泥炭の成因について適切な記述を選ぶ。

\Needspace{4\baselineskip}\paragraph{確認した条件・記述（独自要約）}\mbox{}\par

\noindent\textbf{方向：} 最も適切なもの。\par

\noindent\textbf{対象：} 土の一般的な名称と成因。\par

\paragraph{問われる核心} 特殊土の母材・成因・残存植物組織

\paragraph{必要知識} 特殊土の母材・成因・残存植物組織。母材・運搬堆積・有機物起源、風化、弱面、水の作用、岩盤と土の物性。

\paragraph{考え方・正答を選ぶ理由（査読待ち初稿）} 色や粒状の見た目だけでなく母材と形成過程を対応させる。まさ土は花崗岩風化、しらすは火山噴出物、黄土は風成シルト、火山灰質粘性土は火山灰の風化等、泥炭は植物遺体の不完全分解である。適切なのは泥炭の5である。

\Needspace{4\baselineskip}\paragraph{選択肢の検討（査読待ち初稿）}\mbox{}\par

\noindent\textbf{1：不適切} まさ土の母材は流紋岩である。。まさ土は主に花崗岩類の風化により砂状となった土である。残留と崩積はその後の移動の区別であり、流紋岩を母材とする記述が誤り。\par

\noindent\textbf{2：不適切} しらすは花崗岩の風化物とその二次堆積物である。。しらすは主に珪長質火砕流の未溶結部などに由来する火山性堆積物で、二次堆積物もある。花崗岩の風化土とは起源が違う。\par

\noindent\textbf{3：不適切} 黄土は海洋堆積物である。。黄土（レス）は風で運ばれたシルトを主体とする風成堆積物で、海成という説明が誤り。\par

\noindent\textbf{4：不適切} 火山灰質粘性土をレスと呼ぶ。。火山灰由来の粘性土と風成シルトのレスを同一にしない。分布地域が一部重なるだけでは成因を一致させられない。\par

\noindent\textbf{5：適切} 泥炭はピートとも呼ばれ、未分解の植物組織を肉眼で判定できる。。水分の多い環境などで植物遺体の分解が遅れ、繊維・組織を残す有機質土が泥炭である。分解程度には幅があり、あらゆる泥炭の全組織が必ず肉眼で見えるという一般化は避ける。\par

\paragraph{再利用する考え方} 特殊土の母材・成因・残存植物組織を成因と観察される性質に結び付け、例示と全称命題、主要鉱物と副成分、公式解答と未解決の解釈を区別する。

\paragraph{共通理論} \hyperref[theory:engineering:soil]{第\ref*{theory:engineering:soil}節}

\par 参考：\cite{mlit_masa,maff_shiras2017,mlit_peat,gsj_loess2002}。

\paragraph{関連問題} \hyperref[q:h27-reg-applied-29]{H27-REG-APPLIED-29}

\par 難易度 3（校準初稿）、優先度未評価。

\subsection{平成27年度 問29}
\label{q:h27-reg-applied-29}

\noindent\textbf{原典：} sources/H27/17\_applied\_science.pdf、PDF物理11頁。

\textbf{公式正答：} 5（single）。

\paragraph{問題の要旨} まさ土・しらす・黄土・火山灰質粘性土・泥炭の成因について適切な記述を選ぶ。

\Needspace{4\baselineskip}\paragraph{確認した条件・記述（独自要約）}\mbox{}\par

\noindent\textbf{方向：} 最も適切なもの。\par

\noindent\textbf{対象：} 土の一般的な名称と成因。\par

\paragraph{問われる核心} 特殊土の母材・成因・残存植物組織

\paragraph{必要知識} 特殊土の母材・成因・残存植物組織。母材・運搬堆積・有機物起源、風化、弱面、水の作用、岩盤と土の物性。

\paragraph{考え方・正答を選ぶ理由（査読待ち初稿）} 色や粒状の見た目だけでなく母材と形成過程を対応させる。まさ土は花崗岩風化、しらすは火山噴出物、黄土は風成シルト、火山灰質粘性土は火山灰の風化等、泥炭は植物遺体の不完全分解である。適切なのは泥炭の5である。

\Needspace{4\baselineskip}\paragraph{選択肢の検討（査読待ち初稿）}\mbox{}\par

\noindent\textbf{1：不適切} まさ土の母材は流紋岩である。。まさ土は主に花崗岩類の風化により砂状となった土である。残留と崩積はその後の移動の区別であり、流紋岩を母材とする記述が誤り。\par

\noindent\textbf{2：不適切} しらすは花崗岩の風化物とその二次堆積物である。。しらすは主に珪長質火砕流の未溶結部などに由来する火山性堆積物で、二次堆積物もある。花崗岩の風化土とは起源が違う。\par

\noindent\textbf{3：不適切} 黄土は海洋堆積物である。。黄土（レス）は風で運ばれたシルトを主体とする風成堆積物で、海成という説明が誤り。\par

\noindent\textbf{4：不適切} 火山灰質粘性土をレスと呼ぶ。。火山灰由来の粘性土と風成シルトのレスを同一にしない。分布地域が一部重なるだけでは成因を一致させられない。\par

\noindent\textbf{5：適切} 泥炭はピートとも呼ばれ、未分解の植物組織を肉眼で判定できる。。水分の多い環境などで植物遺体の分解が遅れ、繊維・組織を残す有機質土が泥炭である。分解程度には幅があり、あらゆる泥炭の全組織が必ず肉眼で見えるという一般化は避ける。\par

\paragraph{再利用する考え方} 特殊土の母材・成因・残存植物組織を成因と観察される性質に結び付け、例示と全称命題、主要鉱物と副成分、公式解答と未解決の解釈を区別する。

\paragraph{共通理論} \hyperref[theory:engineering:soil]{第\ref*{theory:engineering:soil}節}

\par 参考：\cite{mlit_masa,maff_shiras2017,mlit_peat,gsj_loess2002}。

\paragraph{関連問題} \hyperref[q:h25-reg-applied-19]{H25-REG-APPLIED-19}, \hyperref[q:r01-ret-applied-28]{R01-RET-APPLIED-28}

\par 難易度 3（校準初稿）、優先度未評価。

\subsection{令和元年度再試験 問28}
\label{q:r01-ret-applied-28}

\noindent\textbf{原典：} sources/R01\_retest/17\_applied\_science.pdf、PDF物理9頁。

\textbf{公式正答：} 3（single）。

\paragraph{問題の要旨} まさ土・しらす・黄土・火山灰質粘性土・泥炭の成因について適切な記述を選ぶ。

\Needspace{4\baselineskip}\paragraph{確認した条件・記述（独自要約）}\mbox{}\par

\noindent\textbf{方向：} 最も適切なもの。\par

\noindent\textbf{対象：} 土の一般的な名称と成因。\par

\paragraph{問われる核心} 特殊土の母材・成因・残存植物組織

\paragraph{必要知識} 特殊土の母材・成因・残存植物組織。母材・運搬堆積・有機物起源、風化、弱面、水の作用、岩盤と土の物性。

\paragraph{考え方・正答を選ぶ理由（査読待ち初稿）} 色や粒状の見た目だけでなく母材と形成過程を対応させる。まさ土は花崗岩風化、しらすは火山噴出物、黄土は風成シルト、火山灰質粘性土は火山灰の風化等、泥炭は植物遺体の不完全分解である。適切なのは泥炭の3である。

\Needspace{4\baselineskip}\paragraph{選択肢の検討（査読待ち初稿）}\mbox{}\par

\noindent\textbf{1：不適切} 黄土は海洋堆積物である。。黄土（レス）は風で運ばれたシルトを主体とする風成堆積物で、海成という説明が誤り。\par

\noindent\textbf{2：不適切} しらすは花崗岩の風化物とその二次堆積物である。。しらすは主に珪長質火砕流の未溶結部などに由来する火山性堆積物で、二次堆積物もある。花崗岩の風化土とは起源が違う。\par

\noindent\textbf{3：適切} 泥炭はピートとも呼ばれ、未分解の植物組織を肉眼で判定できる。。水分の多い環境などで植物遺体の分解が遅れ、繊維・組織を残す有機質土が泥炭である。分解程度には幅があり、あらゆる泥炭の全組織が必ず肉眼で見えるという一般化は避ける。\par

\noindent\textbf{4：不適切} まさ土の母材は流紋岩である。。まさ土は主に花崗岩類の風化により砂状となった土である。残留と崩積はその後の移動の区別であり、流紋岩を母材とする記述が誤り。\par

\noindent\textbf{5：不適切} 火山灰質粘性土をレスと呼ぶ。。火山灰由来の粘性土と風成シルトのレスを同一にしない。分布地域が一部重なるだけでは成因を一致させられない。\par

\paragraph{再利用する考え方} 特殊土の母材・成因・残存植物組織を成因と観察される性質に結び付け、例示と全称命題、主要鉱物と副成分、公式解答と未解決の解釈を区別する。

\paragraph{共通理論} \hyperref[theory:engineering:soil]{第\ref*{theory:engineering:soil}節}

\par 参考：\cite{mlit_masa,maff_shiras2017,mlit_peat,gsj_loess2002}。

\paragraph{関連問題} \hyperref[q:h27-reg-applied-29]{H27-REG-APPLIED-29}

\par 難易度 3（校準初稿）、優先度未評価。

